\chapter{World Regions and Cultures}
\label{ch:world-regions}

The world of \textit{Fate's Edge} is a tapestry of ancient empires, emerging kingdoms, and untamed wilderness. This chapter surveys major regions and cultures that shape the setting---from the marble cities of Ecktoria to the mist-shrouded fields of Aelinnel. These frameworks are yours to adapt, blend, or reimagine.

\section{The Amaranthine Inland Sea}
\index{Amaranthine Coastway}\index{inland sea}

At the heart of the known world lies the \textbf{Amaranthine Inland Sea}, a wind-gnarled waterway ringed by marble quays, vineyard hills, and smoke-blue mountains. For millennia it has served as the circulatory system of trade, faith, and conquest. Tides are subtle, but seasonal winds and river-feeds set the rhythm of commerce, pilgrimage, and war.

\section{Northern Shore of the Amaranthine Sea}

\subsection*{Ecktoria — The Utaran Imperium Successor}
\index{Ecktoria}

Once the furnace of empire (\emph{Marble \& Fire}), Ecktoria remains a palimpsest of power: old stones bearing new banners, old laws written under fresh seals. Though imperial reach waned, its civic habits endure — a echo of late Roman/Byzantine grandeur, with provincial governors, standing bureaucracies, and ever-present legal codes.

\begin{description}
\item[Marble Cities] Forums, amphitheaters, and aqueducts yet flow. District fountains double as public oaths guaranteed by guild charters.
\item[Imperial Roads] Mile-markers of white granite, way-shrines and customary tolls noted for couriers of the \emph{Ashen Staves}.
\item[Legal Legacy] The \textbf{Utaran Civic Codes} govern contracts, inheritance, and war-rights; local custom bends them under licensed \emph{variance}.
\item[Architectural Wonders] Sun-bridges spanning deltas, the \textbf{Vault of a Thousand Maps}, and the \textbf{Amber Arch} petrified by alchemical storm.
\end{description}

\subsection*{Acasia — "The Broken Province"}
\index{Acasia}

Fallen Province. Frontiers braided from roads, rivers, and resentments — a land of petty kingdoms and free companies, reminiscent of the Welsh Marches or the post-Roman frontier. Here the outer seams of empire frayed first. Fortresses turned manors, manors turned townholds, and banners multiplied like thistles after rain.

\begin{description}
\item[Petty Kingdoms] Dozens of river-vales ruled by river-kings and banner-queens. Alliances shift with marriages, harvests, and omens.
\item[Fortified Towns] Walls for defense, not display. Gate-streets kink for ambush; towers carry horn-codes every child knows.
\item[Mercenary Culture] Free companies keep a \emph{Black Ledger}: contracts fulfilled, oaths kept, debts paid.
\item[Cultural Mix] Imperial rites meet clan feasts; old gods share niches with civic saints. Exiles and second chances (\emph{see} Silkstrand tales) are common.
\end{description}

\subsection*{Vhasia — "Old Vhasia \& The Bloodlands"}
\index{Vhasia}

Politically fractured kingdom of courtly intrigue and martial tradition, where ancient bloodlines vie for supremacy amid shifting alliances and ceremonial warfare. Travelers from across the Amaranthine often note its chivalric pageantry, its elaborate dueling codes, and its love of rhetoric — echoes of a certain western kingdom across the channel.

\begin{description}
\item[Fortress Castles] Stone keeps crowned with gilded spires; courtiers plot in tapestried halls while knights train in courtyards.
\item[Political Intrigue] Complex web of alliances, vendettas, and ceremonial duels that settle matters of honor and succession.
\item[Court Culture] Elaborate ceremonies, patronage of arts, and rigid social hierarchies maintained through ritual and reputation.
\item[Heraldic Traditions] Complex system of banners, titles, and precedence that govern social interactions and military commands.
\end{description}

\subsection*{Thepyrgos}
\index{Thepyrgos}

Province and capital city renowned as a center of learning, magic, and scholarly pursuit, where ancient towers house both wisdom and dangerous secrets. A blend of Byzantine bureaucracy and Mediterranean university culture — the City of a Thousand Stairs.

\begin{description}
\item[Scholarly Traditions] Tower-cities where mages, philosophers, and researchers pursue knowledge in specialized colleges and scriptoriums.
\item[Arcane Heritage] Deep traditions of magical study, with libraries containing texts predating the fall of ancient empires.
\item[Academic Rivalries] Intense competition between schools of thought, often manifesting in formal debates, magical duels, or scholarly contests.
\item[Mystical Dangers] Forbidden knowledge and experimental magic that sometimes escape control, creating ongoing threats.
\end{description}

\subsection*{Viterra — "The Last Kingdom"}
\index{Viterra}

Realm that straddles the Dolmis and Amaranthine seas, known for its legalistic approach to governance and strategic river crossings. Its common law — the \emph{Hedge-Law} — and its tradition of legal precedent evoke the courts of a certain misty island kingdom.

\begin{description}
\item[Hedge-Law Culture] Complex system of legal precedents, tolls, and river rights that govern everything from trade to personal conduct.
\item[Duchy System] Semi-autonomous regions governed by dukes who maintain their own courts and armies while owing fealty to the crown.
\item[River Commerce] Economy built around controlling strategic crossings, ferry rights, and maritime trade routes.
\item[Legalistic Politics] Intrigue centered on court cases, charter disputes, and the interpretation of ancient laws rather than open warfare.
\end{description}

\subsection*{Ubral — "The Stone Between Spears"}
\index{Ubral}

Highland realm of rugged clans and fortified holds, where honor culture and martial traditions dominate social interactions — a landscape of braes and glens, of cairns and clan feuds.

\begin{description}
\item[Clan Strongholds] Fortified positions in mountain passes and high valleys, each clan maintaining its own laws and customs.
\item[Honor Culture] Society built around concepts of personal honor, family reputation, and the resolution of disputes through formal challenges.
\item[Highland Warfare] Military traditions emphasizing heavy infantry, defensive positions, and knowledge of mountain terrain.
\item[Clan Loyalties] Complex web of alliances, blood-feuds, and marriage pacts that shift with each generation.
\end{description}

\subsection*{Kahfagia — "The Empire of Wakes and Storm-Flags"}
\index{Kahfagia}

Maritime empire built on naval supremacy and exploration, where ship captains and merchant-adventurers shape both policy and culture — a Venetian–Portuguese hybrid, with lantern‑law and pilot guilds ruling the waves.

\begin{description}
\item[Naval Supremacy] Military and economic power based on controlling sea lanes, harbors, and maritime trade routes.
\item[Explorer Culture] Tradition of venturing into unknown waters, mapping new territories, and establishing trading posts.
\item[Storm-Flag Protocol] Complex system of maritime signals, weather prediction, and naval customs that govern seaborne activities.
\item[Mixed Heritage] Cosmopolitan society influenced by contacts with distant lands and diverse cultures encountered through exploration.
\end{description}

\subsection*{Theona — "The Marsh Crown"}
\index{Theona}

Three island realms connected by causeways and maritime traditions, where wetland resources and naval culture define daily life. The grey skies, peat bogs, and ringforts of Theona echo the landscapes of ancient Éire — a land of high kings, saints, and folk horrors lurking beneath the mist.

\begin{description}
\item[Marsh Agriculture] Sophisticated systems of dikes, canals, and floating gardens that support dense populations in wetland environments.
\item[Island Culture] Distinct traditions for each island, unified by shared maritime customs and inter-island trade.
\item[Waterborne Commerce] Economy based on fishing, water transport, and control of strategic waterways between islands.
\item[Folk Horror Traditions] Deep connection to marsh spirits, water deities, and ancient practices that blur the line between protection and appeasement.
\end{description}

\subsection*{The Mistlands — "Fields Under a Moving Sky"}
\index{Mistlands}

Isolated region shrouded in perpetual mists, where ancient Aelerian protectorate status creates tension between autonomy and oversight. The Mistlands draw from Celtic and Breton folklore — a land of bells, standing stones, and the ever‑present threat of what walks in the fog.

\begin{description}
\item[Mistbound Geography] Landscape of bogs, waterways, and hidden settlements connected by causeways and boat paths.
\item[Bell Culture] Complex system of bells and wards must be maintained to keep the Direwood horrors at bay.
\item[Isolation Tensions] Cultural friction between desire for independence and practical need for trade and protection.
\item[Ancient Secrets] Ruins and artifacts predating the Aelerian protectorate, hinting at older civilizations and forgotten magics.
\end{description}

\section{Peoples and Cultures}

The peoples of Fate’s Edge are many. Some trace their lineage to the earliest days of the Utaran Imperium; others arrived across the sea or emerged from the mist‑hung valleys of the Direwood. What follows are the major cultural groupings — self‑named and deeply particular. Outsiders may use simpler labels; the wise learn the names that people call themselves.

\subsection*{The Lethai — The Sundered Kin}
\index{Lethai}\index{Elves}

The Lethai are the eldest of the mortal lineages, their memory stretching back before the fall of empires now forgotten. An ancient constraint divides them: no Lethai may bear both the Gift of the Body and the Gift of the Mind. From this sundering grew three distinct peoples, each with its own philosophy, homeland, and customs.

\subsubsection*{Lethai-al — “People of the Body” (Wood Elves)}
\index{Wood Elves}\index{Lethai-al}

The Lethai-al dwell in the **Valewood**, that ancient forest where roof‑trees braid the sky and rivers think aloud. Their memory is arboreal — rings, seasons, coppice ledgers — and their oaths are \emph{root‑law}: debts counted in years, paid in living work. To outsiders they seem quiet; to neighbors they are relentless auditors of footprint and flow.

\begin{description}
\item[Forest Harmony] Lifestyle integrated with woodland ecosystems, practicing sustainable hunting, gathering, and cultivation. They do not “use” the forest; they live within its clauses.
\item[Body‑Centric Philosophy] Emphasis on physical experience, instinct, and the wisdom of sinew and breath over abstract thought. A Lethai‑al proverb: \emph{“The hand that fells the tree must plant the sapling.”}
\item[Living Magic] Spellcasting works \emph{with} natural forces, not against them — plant‑speech, animal kinship, and the slow shaping of living wood.
\item[Seasonal Rituals] A calendar of ceremonies marking planting, harvest, and the winter quiet. Outsiders are welcome at the Autumn Trade Moot, but the Spring Rites are veiled.
\item[Homeland] The **Valewood** itself — no single capital, but shifting grove‑courts and deep hollows where the oldest trees remember empire.
\end{description}

\subsubsection*{Lethai-thora — “People of the Mind” (High Elves)}
\index{High Elves}\index{Lethai-thora}

The Lethai‑thora make their homes in cities — chiefly **Thepyrgos**, the City of a Thousand Stairs — where they weigh arguments like bridges and translate other peoples’ law into forms that carry. Their courts count consequences; their speech is context‑saturated and exact.

\begin{description}
\item[Scholarly Excellence] Deep traditions of academic study, magical research, and philosophical debate. The Great Library of Thepyrgos holds texts that predate the Sundering.
\item[Mind‑Centric Philosophy] Reason, logic, and the pursuit of abstract knowledge over physical concerns. A Lethai‑thora axiom: \emph{“A clean proof is worth a thousand swords.”}
\item[Arcane Mastery] Advanced magical techniques and theoretical understanding that often surpass other traditions. Many of the Free Casters’ foundational texts were penned by Lethai‑thora hands.
\item[Long Perspective] Tendency to view problems and conflicts through the lens of centuries rather than immediate gain. This patience can be wisdom — or paralysis.
\item[Homeland] **Thepyrgos** and its satellite academies, with enclaves in other major cities where scribal work and jurisprudence are valued.
\end{description}

\subsubsection*{Lethai-ar — “The Oathbound” (Dark Elves)}
\index{Dark Elves}\index{Lethai-ar}

The Lethai‑ar are not a separate bloodline but a vow‑bound cadre. They are Lethai‑al or Lethai‑thora who step off those paths to live under threshold patrons — **Inaea** the Weaver and **Isoka** the Serpent. In some ages they are scarce; in others — like the present — they gather wherever borders fray and oaths need teeth. They are known for \emph{mask‑right}, precise courtesies, and the unsettling efficiency of vows that bind places, routes, and roles.

\begin{description}
\item[Serpent Wisdom] Followers of Isoka emphasize transformation, renewal, and the shedding of old identities for new growth. A shedding can be a cure, an escape, or a cruelty — context is all.
\item[Spider Webs] Devotees of Inaea focus on connections, patterns, and the weaving of fate through careful manipulation. Mercy with memory; knots that mend before they bind.
\item[Philosophical Balance] Neither inherently good nor evil, but representing alternative approaches to power and influence. A Lethai‑ar may be a healer or a poisoner, a mediator or a cut‑throat — the method, not the morality, defines them.
\item[Cultural Rarity] Uncommon in most settings, often viewed with suspicion or fascination. Their masks and vows unsettle those who prefer straightforward dealings.
\item[Homeland] No fixed homeland. They gather in **mask‑circles** within other Lethai communities or in hidden enclaves near ancient thresholds.
\end{description}

\subsection*{The Aeler — People of Stone}
\index{Dwarves}\index{Aeler}

The Aeler are the mountain‑born, built like their peaks — layered, load‑bearing, enduring. They call their home **The Mountains**, **Undermount**, or **The Stonelands**; outsiders often say “Aelerian range,” but that is a cartographer’s convenience, not a name the Aeler use for themselves.

\begin{description}
\item[Stone‑Sense] Innate ability to understand and work with geological formations, making them master miners, architects, and engineers. An Aeler can read a tunnel’s stress like a human reads a face.
\item[Clan System] Complex social structure based on family lines, with intricate systems of honor, debt, and mutual obligation. Every clan keeps a **Keystone Ledger** — a record of favors owed and paid.
\item[Craft Traditions] Legendary skills in metalworking, stonework, and engineering refined over generations. Aeler steel is sought across the Amaranthine; Aeler bridges are the last to fall in a siege.
\item[Underground Cities] Vast complexes carved from mountain hearts, connected by tunnels and halls that serve both practical and ceremonial purposes. The city of **Khaz‑Vurim** is the oldest and greatest.
\item[Deep Law] Contracts are carved, not inked. A broken oath can be read in the stone for centuries. Justice is a matter of measurement — a tilted beam, a mis‑weighted coin.
\end{description}

\subsection*{The Aelinnel — People of Sums}
\index{Gnomes}\index{Aelinnel}

The Aelinnel dwell in **Aevrossa**, the hawthorn hills south of the Valewood — a land of living wood, worked stone, and air that tastes of mathematics. Their lives intertwine with **courtesy, copper, and count**.

\begin{description}
\item[Mathematical Culture] Society built around complex calculations, probability, and the belief that all phenomena can be understood through numerical relationships. A bridge’s strength, a market’s price, a court’s verdict — all are sums.
\item[Fey Logic] Non‑linear thinking patterns that seem illogical to outsiders but follow their own internal consistency. \emph{“If you speak the ninth name, you must also speak the silence that follows it.”}
\item[Mist Adaptation] Unique abilities to navigate and manipulate the perpetual mists of Aevrossa. They read fog like a ledger, finding paths where others see only white.
\item[Contract Culture] Deep tradition of precise agreements, wordplay, and the careful crafting of obligations. Every bargain has two ledgers: what was \emph{said} and what was \emph{meant}.
\item[Homeland] **Aevrossa** — a realm of stone circles, antler‑posts, and tide‑cut stairs. Its capital is the **Hall of Aelinnel**, a timber keep threaded between standing stones.
\end{description}

\subsection*{The Aelaerem — People of the Hearth}
\index{Halflings}\index{Aelaerem}

The Aelaerem are a people of movement and assembly, living among gentle slopes and hedged lanes. Their homeland’s true name is **hushed** — they refer to it only with euphemisms, for to speak its old name is to invite the attention of the Neighbors. Outsiders may call it “the Halfling Downs” or “the Southern Valleys”; the Aelaerem smile and say nothing.

\subsubsection*{Euphemisms for the Homeland}
The Aelaerem use many pseudonyms, each a quiet oxymoron that hides a darker truth:

\begin{center}
\begin{longtable}{|p{3cm}|p{5cm}|p{5cm}|}
\hline
\textbf{Common Name} & \textbf{Hidden Meaning} & \textbf{Old Tongue} \\
\midrule
The Silent Welcome & “Welcome” = hospitality; “Silent” = calm, a gentle invitation. Sounds like a village festival. & \emph{Rath‑Gor} (“the binding mouth”) — the ancient seal where the Neighbors are held. \\
The Merry Maw & “Merry” = good cheer, feasting; “Maw” = a large gaping mouth. Sounds like a tavern name. & \emph{Gor‑Elen} (“the sealed mouth”) — the cursed cavern where eldritch doors are locked. \\
Stillwater (or The Still‑Water) & “Still” = peace; “water” = life, farming. A common halfling landmark. & \emph{Nok‑Lur} (“the black source”) — an underground well that is actually a conduit to the other world. \\
The Quiet Hearth & “Hearth” = home, warmth; “Quiet” = serene, safe. & \emph{Ael‑Gren} (“the sealed fire”) — the ancient stone ring that locks the Neighbors away. \\
The Gentle Fold & “Fold” = a shepherd’s safe enclosure; “Gentle” = pastoral. & \emph{Thar‑Gur} (“the gnashing cleft”) — the fissure in the earth that serves as the “mouth” of the seal. \\
\hline
\end{longtable}
\end{center}

\noindent The Aelaerem themselves are warm, generous, and fiercely protective of their own — but their land hides a hunger. Hospitality is their public law; beneath it runs an older hedge‑law of cup‑marks, red thread, and the quiet attention of the Neighbors.

\begin{description}
\item[Hearth Culture] Deep connection to home, family, and the maintenance of traditional ways. A red door promises bread, salt, and one safe night if you come honest.
\item[Agricultural Expertise] Sophisticated farming techniques and seasonal celebrations that mark the rhythm of rural life. Their cider, wool, and beeswax are known across the Amaranthine.
\item[Folk Horror Elements] Dark undercurrents in seemingly peaceful communities, where hospitality can become a trap and tradition carries hidden costs. The Neighbors are always watching.
\item[Community Defense] Strong traditions of mutual aid and collective action when the community is threatened. When the bells ring, every Aelaerem takes up a tool or a weapon.
\end{description}

\subsection*{Other Peoples — “Beyond the Old Roads”}
\index{Other Races}

The lands beyond the Amaranthine and the Aelerian passes hold many other folk. Some are recent arrivals; others are ancient survivors of forgotten cataclysms.

\begin{description}
\item[Nomadic Tribes] Various peoples who follow seasonal patterns across steppes, deserts, and other marginal lands. The **Ykrul** are the most numerous, but they are far from alone.
\item[Coastal Peoples] Maritime cultures that live in harmony with ocean environments and maintain their own naval traditions. The **Linn** and **Zakov** corsairs, the **Kahfagian** pilots, and the island clans of the **Dolmis**.
\item[Border Cultures] Mixed communities that arise where different major cultures meet, creating unique hybrid traditions. The **Vilikari** straddle empire and steppe; the **Acasian** marches brew their own fierce independence.
\item[Ancient Survivors] Remnants of older civilizations that persist in isolated regions, maintaining forgotten knowledge and customs. The **Mistlands** bell‑wardens, the **Theonan** marsh‑kings, and the silent watchers of the **Direwood**.
\end{description}
\section{Regional Specialties and Resources}

\subsection*{Climate and Biomes of the Amaranthine}

The Amaranthine Inland Sea is ringed by fertile coasts, rolling hills, and mountain ranges. Unlike the arid southern shores of Old Earth’s middle sea, the wind patterns here reverse — the hot, dry air of the Sahel belt is pushed south, leaving the northern and eastern shores with ample rainfall and the southern shores as temperate grasslands. Thus, the lands around the Amaranthine are lush and productive, from the olive groves of Ecktoria to the wheat fields of Acasia to the horse pastures of the Vilikari steppe.

\subsection*{Economic Strengths by Region}

\begin{description}
\item[Ecktoria] \textbf{Marble, legal services, grain, olive oil, wine, purple dye (murex), glassworks, coinage.} Ecktoria’s marble quarries supply the entire sea. Its law courts and coin‑houses handle contracts for half the continent. The purple dye, extracted from sea snails along the Sun Coast, is worth its weight in silver.

\item[Acasia] \textbf{Cattle, horses, wool, iron, salt, mercenary contracts, border tolls.} The Broken Marches are poor in coin but rich in livestock and metal. Its free companies are the finest light cavalry in the known world. The salt pans of the western marches supply the interior.

\item[Vhasia] \textbf{Wine, timber, stone, cloth (linen and wool), illuminated manuscripts, tapestries.} The river valleys of Vhasia produce excellent vintages; its forests provide oak for shipbuilding. The scriptoria of its abbeys are famed for their calligraphy and illuminated manuscripts.

\item[Thepyrgos] \textbf{Scholarly texts, magical reagents, glassware, precision instruments, silk (imported and finished), philosophical schools.} The City of a Thousand Stairs is a hub of learning; its universities license thaumaturgical research. The glassblowers of Thepyrgos produce lenses and astrolabes unmatched anywhere.

\item[Viterra] \textbf{Tin, copper, wool, tin‑glazed pottery, legal charters, insurance contracts, river tolls.} The river crossings of Viterra control trade between the inland sea and the northern ocean. Its duchy courts are famous for settling disputes without bloodshed — usually.

\item[Ubral] \textbf{Lead, silver, slate, hardy mountain sheep (wool and cheese), whisky, mercenary mountaineers.} The high clans of Ubral mine the peaks for silver and lead; their whisky is legend. Their clansmen are sought as scouts and skirmishers in mountain warfare.

\item[Kahfagia] \textbf{Shipbuilding, naval stores (pitch, tar, sailcloth), spices (imported), exotic woods, dyes (indigo, cochineal), potatoes, maize, tomatoes.} The Kahfagian fleets control the western sea lanes. Their merchants bring goods from across the “west lands” — new world crops that have revolutionized agriculture in recent centuries.

\item[Linn] \textbf{Fish (dried and salted), whale oil, amber, walrus ivory, furs, raiding, ship‑tolls.} The Linn skerries are rich in marine resources. Their longships carry furs from the northern forests and amber from the Baltic coasts.

\item[Theona] \textbf{Citrus fruits, olive oil, pottery (amphorae), salt, peat, glass beads, star‑compasses.} The three islands of Theona produce citrus oils and fine ceramics. Their star‑compasses, crafted in the cliffside observatories, are essential for deep‑sea navigation.

\item[Mistlands] \textbf{Peat, bog iron, fog‑pollen ink, mistglass (humming glass), woven fog‑wool, herbal remedies (nightshade, feverfew, belladonna).} The Mistlands are resource‑poor but materially unique. Mistglass — glass that sings in storms — is prized by sea captains. Fog‑wool, harvested from mist‑spirits, is impossibly light and warm.

\item[Zakov] \textbf{Ship salvage, black market goods, exotic poisons, reef coral, pearls, smuggled luxuries.} The pirate‑city of Zakov lives on stolen cargo and “found” wrecks. Its coral reefs yield pearls and rare sea silks. No question too dark, no cargo too illicit — for the right price.

\item[Vilikari (Frontier Marches)] \textbf{Horses, leather goods, horn, bone, honey, wax, mercenary light cavalry.} The Vilikari straddle the steppe and the sown. Their horses are swift and hardy; their honey and wax are traded across the marches. Vilikari light horse are the eyes of any army.

\item[Ykrul (Steppe Nomads)] \textbf{Horses (the finest), hides, wool felt, dairy products (cheese, yogurt), dried meat, amber, captive trade (slaves, though the Ykrul themselves dislike the term and practice limited bondage).} The Ykrul herds are the wealth of the Violet Steppe. Their horses are exported as far as Silkstrand. Amber found in river deposits is carved into amulets and beads.

\end{description}

\subsection*{Unique Products (Expanded)}

\begin{itemize}
\item \textbf{Ecktoria:} \emph{Purple marble} (only from the Imperial Quarry), \emph{ceramic plumbing} (self‑cleaning clay pipes), \emph{state bonds} (negotiable instruments guaranteed by the Ecclesiarchy of the Everflame).
\item \textbf{Acasia:} \emph{Curse‑iron} (cold iron from the Blackwood, which drips rusty water when near fae or undead), \emph{hedge‑knight contracts} (bought and sold like commodities on the exchange at Silkstrand).
\item \textbf{Vhasia:} \emph{Illuminated bestiaries} (some of which move or speak), \emph{vengeance wine} (fermented with grape leaves from a murdered man’s vineyard — said to carry his spirit).
\item \textbf{Thepyrgos:} \emph{Echo glass} (panes that repeat the last words spoken before them), \emph{thousand‑step tea} (a brew that sharpens the mind for exactly one thousand heartbeats).
\item \textbf{Viterra:} \emph{Hedge‑law scrolls} (legal documents that quietly update their clauses as circumstances change), \emph{river silver} (coins minted from electrum found only in the Belworth).
\item \textbf{Ubral:} \emph{Whisky of the Twelve Clans} (aged in lead‑lined casks, gaining a faint, sweet toxicity that drinkers call “the bite”), \emph{slate that echoes} (roofing tiles that repeat the sounds of a house’s history).
\item \textbf{Kahfagia:} \emph{Storm‑flags} (banners that predict weather by changing color), \emph{lantern‑glass} (treated glass that glows for a year after one night’s exposure to the order’s beacons).
\item \textbf{Linn:} \emph{Whale‑bone runes} (divination tools that must be moistened with blood or saltwater to work), \emph{frost silk} (spider silk harvested from ice caves, nearly invisible and unbreakably cold).
\item \textbf{Theona:} \emph{Star‑compasses} (brass instruments that point to the nearest shrine of the Navigator Saint), \emph{amphorae of silence} (clay jars that, when broken, create a zone of magical quiet for one minute).
\item \textbf{Mistlands:} \emph{Fog‑pollen ink} (invisible except under a lantern or when breathed upon), \emph{mistglass humming} (each glass has its own tone; shattering one creates a sound that can be heard from a mile away).
\item \textbf{Zakov:} \emph{Black coral amulets} (ward against drowning), \emph{songs of the reef} (knot‑coded messages that only other pirates can read), \emph{pearl of forgetfulness} (grants one short‑term memory erasure when crushed).
\item \textbf{Vilikari:} \emph{Horse‑hide scrolls} (waterproof and nearly indestructible), \emph{blood‑honey} (mead made with honey from hives that fed on corpse‑flowers; potent and hallucinogenic).
\item \textbf{Ykrul:} \emph{Kon'reh boards} (carved stone or ivory game boards whose pieces move according to mathematical rules), \emph{wake‑ropes} (enchanted ropes that float and can be used to tie two moving objects together mid‑gallop).
\end{itemize}

\subsection*{Trade and Resource Clocks}

Certain resources can be tracked with clocks to create economic pressure:

\begin{itemize}
\item \textbf{Grain Price (4)}: If supply drops, riots in cities; if oversupply, farmers rebel.
\item \textbf{Salt Monopoly (6)}: Control of salt pans gives leverage over an entire region.
\item \textbf{Fleet Readiness (6)}: A naval power’s ability to project force.
\item \textbf{Market Scrip (4)}: Confidence in paper currency; collapse leads to barter.
\end{itemize}

\section{Travel and Trade}

\subsection*{Major Routes (Reference)}
\begin{description}
\item[Amaranthine Coastway] Sea‑corridor linking Theona, Linn, Zakov, Kahfagia, and Ecktoria. Sheltered anchorages every fifty miles. Seasonal winds run clockwise.
\item[Astroegro Straits] The narrows between Thepyrgos and the northern shore. Pilot‑ruled; tolls vary by the tide, the weather, and the pilot’s mood.
\item[River Roads] Grain and scrip from interior markets to sea. The Belworth, the Yloka, and the Acadian rivers are the main arteries.
\item[Aelerian Passes Underways] Vault‑routes beneath the peaks. Sealed in deep winter, but in summer they offer the safest passage east–west.
\item[Kahfagian Sea Lanes] Goods from across the “west lands” — potatoes, maize, cochineal, and stranger wares. Dangerous, but the profits are legendary.
\item[The Way of Silk] Overland routes from the far east, ending in Silkstrand. Caravans take two years to cross the steppe and the mountains.
\item[Shadow Corridors] Liminal shortcuts near the Ways Between. Risky, fast, and never the same twice. A day’s travel may be an hour — or a month.
\end{description}

\subsection*{Travel Considerations}
\begin{itemize}
\item \textbf{Road Quality}: Utaran imperial highways still exist in Ecktoria and parts of Vhasia. Further east, they become gravel, cart‑ruts, or mere tracks. In the Mistlands, the “road” is a causeway of reeds and stone that floods at each tide.
\item \textbf{Bridges \& Ferries}: Every crossing is a toll point. Some bridges are magical — they only open to certain keys or at certain hours. Ferries may be rowed by the living — or by the dead.
\item \textbf{Seasonality}: The Amaranthine is navigable year‑round, but winter storms can close the straits. Mountain passes close from First Snow to Spring Thaw. The steppe is impassable during the spring melt (rasputitsa).
\item \textbf{Safe Havens}: Caravanserais are common along the Way of Silk. In the north, monasteries offer shelter. On the coast, lighthouse‑cloisters owe aid to shipwrecked sailors by ancient charter.
\item \textbf{Customs and Quarantines}: Entering a new region often requires papers — a safe‑conduct, a merchant seal, or a ritual cleansing. The Mistlands may quarantine travelers for three days to ensure they are not possessed by fog‑spirits.
\end{itemize}

\section{Regional Clocks and World Response}
\index{world response}\index{clocks}

Tie adventures to \textbf{clocks} that change the map when they fill. Each clock is a narrative front; when it fills, the region’s situation escalates.

\begin{itemize}
\item \textbf{Grain Shortage (4)}: If filled, food riots reshape a Theonan city’s politics. May spawn a famine if left unchecked.
\item \textbf{Banner‑Muster (6)}: Steppe tribes unify under a Khagan. Caravans demand new terms; raids increase; travel becomes dangerous.
\item \textbf{Harbor Scrip Crash (4)}: Trade letters lose value; smugglers thrive; debts are called in. Banks may fail, and noble houses may fall.
\item \textbf{Fog‑Roused (6)}: Mistland spirits awaken. Bells fail at night; offerings must be increased. The Direwood’s influence spreads.
\item \textbf{Stone‑Press (6)}: Aeler deep mines begin a slow, ticking collapse. Ventilation fails; old tunnels seal; clans feud over responsibility.
\item \textbf{Chevauchée (4)}: Vhasian war‑bands ride. Border villages burn; the king’s justice is distant. Mercenary companies profit.
\item \textbf{Reef Shift (6)}: Zakov’s coral reefs shift, exposing new wrecks — and closing old passages. Salvage rights become a blood sport.
\end{itemize}

\smallskip
\noindent GM Story Beats may \emph{tick} these when the party’s choices strike regional nerves: missed payments, broken oaths, loud magic, conspicuous success, or even a well‑timed rumor.

\section{Cultural Practices and Customs}

\subsection*{Languages and Cant}
\begin{itemize}
\item \textbf{Utaran High} (court, law, scholarship; spoken in Ecktoria, Thepyrgos, and by educated elites across the sea).
\item \textbf{River Cant} (trade pidgin with gesture‑signs; used on the Amaranthine and major rivers).
\item \textbf{Steppe Tongues} (whistled across distance; each clan has its own variations).
\item \textbf{Sea‑Patter} (mariners’ clipped code; used by Linn, Zakov, and Kahfagian sailors).
\item \textbf{Bell‑Speech} (Mistlands only; tones and rhythms of bells convey legal and spiritual messages).
\item \textbf{Hedge‑Law Notation} (Viterra and parts of Vhasia; legal abbreviations written in the margins of charters).
\end{itemize}

\subsection*{Religious Patterns}
\begin{itemize}
\item \textbf{Temple \& School} (Theona): Ethical schools blend with temple tithes; festivals mark navigation seasons. Pilgrims walk causeways at low tide.
\item \textbf{Ancestor Poles} (Vilikari): Kin‑spirits seated at feasts; oaths taken under the watch of names. The poles are carved with lineages stretching back centuries.
\item \textbf{Sky \& Earth} (Steppe): Wind‑knots bound, libations poured into the first hoofprint. The Ykrul revere the Violet Sky and the endless grass.
\item \textbf{Sea‑Rites} (Linn/Zakov): Weather judgments witnessed by storm‑priests. Offerings of salt and silver are thrown overboard before a voyage.
\item \textbf{The Everflame} (Ecktoria, Thepyrgos, parts of Vhasia): Cults of the Eternal Fire, with lamp‑tenders and ash‑priests. Judgment by flame is a common legal ordeal.
\item \textbf{Fog Faiths} (Mistlands): Bells and fog‑names soothe fears; spiritualists mediate with the grey things. The dead are not buried — they are rung away.
\end{itemize}

\subsection*{Law and Custom}
\begin{itemize}
\item \textbf{Utaran Codes}: Contracts, inheritance, civic duties; licensed \emph{variance} lets local custom lean the written law. Common in Ecktoria, Thepyrgos, and former imperial provinces.
\item \textbf{Clan Law} (Acasia, Vilikari, Ubral, Ykrul): Oath‑payment by cattle, steel, or service‑days. Blood price is fixed for each injury; the victim’s clan must accept it or escalate.
\item \textbf{Merchant Law}: Arbitration by bonded factors; ledgers sealed in wax and salt. Defaulting on a debt inscribed in a Factor’s Ledger is a crime in every port.
\item \textbf{Hedge‑Law} (Viterra, parts of Vhasia): Precedents set by parish juries and traveling justices. A hedge decree is binding until the next quarter court.
\item \textbf{Lamp‑Law} (Mistlands, Kahfagian lantern precincts): Rules enforced by the lighting or extinguishing of public lamps. A dark lane is a legal vacuum.
\end{itemize}

\section{Magic and Attitude by Region}
\index{magic!cultures}

\begin{itemize}
\item \textbf{Ecktoria}: Licensed thaumaturges file \emph{Casting Notices} for urban work; unfiled magic draws fines — and Everflame attention. Magic is seen as a civic utility, like plumbing.
\item \textbf{Acasia}: Folk‑wards respected; visible sorcery can start a levy. Hedge witches are tolerated; university mages are viewed with suspicion.
\item \textbf{Aeler}: Rituals fold into craft; backlash treated like a cracked beam — fix it, document it, and pay the guild fine. Magic is a trade like any other.
\item \textbf{Mistlands}: Bells and fog‑names soothe fears; spiritualists mediate with the grey things. Unlicensed magic is punished by servitude to the bell‑wardens.
\item \textbf{Linn / Zakov}: Weather rites must be witnessed by three wind‑bearers; false rites are crimes against the ship. Magic at sea is tightly regulated by the storm‑priests.
\item \textbf{Thepyrgos}: Academics argue about proper magical forms. Unorthodox casting is not illegal — merely unfashionable — and can ruin a scholar’s reputation.
\item \textbf{Vhasia}: The Church of the Flame distrusts arcane magic; secular lords employ court mages anyway. A delicate balance of public condemnation and private patronage.
\item \textbf{Viterra}: Hedge‑magic is common and unregulated; any thaumaturgy that leaves a trace must be registered with the local sheriff. Penalties are fines, not prison.
\item \textbf{Kahfagia}: The Lantern‑Law tolerates magic that assists navigation or commerce. Weather magic is licensed; necromancy is punished by drowning.
\item \textbf{Ubral}: Mountain clans view magic as a weapon, not a tool. Outsider mages are watched; clan shamans are respected but feared. The death penalty for using magic to steal clan property.
\item \textbf{Vilikari}: Mixed attitudes — some clans employ shamans; others see magic as dishonorable. The Federate King has a court astrologer but no official battle mages.
\item \textbf{Ykrul}: The Kon'reh masters treat magic as geometry; precise, bounded, and ethically neutral. Wild magic is distrusted; the Wake‑law forbids raising the dead.
\end{itemize}

\section{Creating Regional Adventures}
\index{regional adventures}

\subsection*{Using Regional Characteristics}
Consider:
\begin{itemize}
\item How geography shapes travel and pursuit (mountains, rivers, passes, seasons).
\item Which customs open doors (or close them) — guest‑right, safe‑conduct, charter.
\item Which tensions tick \textbf{clocks} — feuds, famines, heresies, succession.
\item Which resources (salt, steel, scrip, secrets) define stakes.
\end{itemize}

\subsection*{Mixing Regional Elements}
\begin{itemize}
\item \textbf{Tin vs. Tide}: Aeler miners accuse a port‑city of short‑weighted scales. The conflict escalates from court to blockade.
\item \textbf{Festival of Threads}: Islanders, Steppe riders, and Heartlanders trade rites and insults under watchful magistrates. A murder must be solved without violating the festival’s truce.
\item \textbf{Fog and Fire}: A Mistlands relic surfaces in a Theonan auction; three factions bid with more than coin — and one bids with a curse.
\item \textbf{The Unquiet Map}: An Aeler survey contradicts an old border. Whose truth stands? The stones speak, but they speak in riddles.
\end{itemize}

\begin{tcolorbox}[colback=blue!5!white,colframe=blue!75!black,title=Regional Adventure Seeds,fonttitle=\bfseries]
\textbf{Mistlands Mystery (Clock 4: Bells Go Silent)}
\begin{itemize}
\item Dusk bells falter. Tracks end at a mirror‑pool that shows tomorrow’s sky.
\item \emph{Complications (SP)}: Fog names stolen; a bell‑tree cracks; a reed altar burns cold.
\item \emph{Resolution}: Offer a true name to the pool or re‑hang the broken bell with copper nails.
\end{itemize}

\textbf{Broken Marches Conflict (Clock 6: Valley Claim)}
\begin{itemize}
\item Two lords court your company. One keeps books; one keeps graves tidy. Both want the same mill.
\item \emph{Complications (SP)}: Ambush at witness trees; the Black Ledger demands a tithe; a feud token is planted in your pack.
\item \emph{Resolution}: Settle the dispute by duel, by ledger audit, or by destroying the mill entirely.
\end{itemize}

\textbf{Stone Kingdom Discovery (Clock 6: Archive Wakes)}
\begin{itemize}
\item An echo‑vault opens to song. Memory‑stones argue \emph{with each other} about who owns the ore.
\item \emph{Complications (SP)}: Rival charter; a cracked beam of magic; miners strike for better rites.
\item \emph{Resolution}: Decode the old arguments, forge a consensus, or collapse the gallery and start anew.
\end{itemize}

\textbf{Theona Trade War (Clock 4: Harbor Scrip Collapse)}
\begin{itemize}
\item Watermarks sing wrong. A counterfeit choir spreads false coin across three islands.
\item \emph{Complications (SP)}: Dock riot; lighthouse shutters locked; a patron’s legal \emph{variance} revoked.
\item \emph{Resolution}: Find the counterfeit master, destroy the press, or redeem the scrip with a cargo of star‑compasses.
\end{itemize}
\end{tcolorbox}

\section{Adapting Regions to Your Campaign}
\index{campaign adaptation}

These regions are \emph{frameworks}. Rename, splice, or tilt them toward your tale:
\begin{itemize}
\item Merge Mistland bells with island wind‑courts to create \emph{storm‑bell law}.
\item Let the Merchants’ Concord adopt steppe arbitration; settle contracts at a gallop.
\item Tie your party to a \textbf{regional clock} — when it fills, the map, and your story, change.
\item Swap resources: maybe the Mistlands have amber instead of peat; maybe Acasia exports tin.
\item Introduce a new region — the “Inner Lake” — by borrowing from the cultures above.
\end{itemize}

Remember: the world should \emph{answer} the players. Roads reroute around their deeds; bells ring differently after they pass; ledgers carry their names in salt. The resources they seize, the trade they disrupt, the customs they honor or break — all of it echoes.

\clearpage
\section{Aeler — People of Stone, Breath, and Ledger}

\subsection*{Background: The Mount-Born Engineers}
Built like their mountains—layered, load-bearing, and enduring—the Aeler are masters of infrastructure, subtle influence, and the deep mathematics of survival. In their underground holds, a lantern's hue is a balance sheet; a bell-note means more than a shout. They call this discipline \textbf{deep accounting}: air tallied by vent-shafts, lamp-time written in chalk, calories measured in the language of ovens.

Above ground, their influence is felt but not always seen. Their work—bridges that do not fail, levees that answer one key, ovens that feed thousands—acts as the hidden bones of cities. They do not conquer through banners, but through the gate that must be paid, the bridge that “politely rests” if its clause is broken, and the surety of grain that keeps a quarter from starving.

\paragraph{Key Cultural Concepts}
\begin{itemize}
  \item \textbf{Tally-Law}: If it isn't written, it isn't owed; if it cannot bear weight, it isn't promised.
  \item \textbf{Keystone Rights}: Control the piece that holds the whole. Maintenance for access; repair for rate.
  \item \textbf{Grain Surety}: Winter ovens and storage domes under public charter; price courts on marked days.
  \item \textbf{Water \& Flood}: Sluice-math is power. They lease keys, not walls; a city opens for trade faster than an army can take it.
  \item \textbf{Mint \& Measure}: Calibrated weights and indelible dies. When coins bear Aeler marks, courts sharpen.
\end{itemize}

\subsection*{Racial Skill Increase}
Choose one:
\begin{enumerate}
  \item \textbf{Stone \& Breath}: +1 die to \textbf{Craft}, \textbf{Tinker}, and \textbf{Survival} when dealing with infrastructure, construction, or resource management. In underground or dense-urban environments, gain \textbf{Position +1} to navigate, maintain, or sabotage systems.
  \item \textbf{Deep Accounting}: Once/scene, spend \textbf{1 Boon} to \emph{audit} a situation—ask one question about hidden resources, costs, or dependencies; the GM answers truthfully.
\end{enumerate}

\subsection*{Thematic Attribute}
Increase either \textbf{Body} or \textbf{Wits} by 1 (to a max of 5). \textbf{Spirit} and \textbf{Presence} unchanged.

\subsection*{Talent: Vent Prior's Training (3 XP)}
\textit{Req: \textbf{Craft} 1+, \textbf{Wits} 2+}
\begin{itemize}
  \item +1 die on checks involving air quality, ventilation, structural integrity, or underground navigation.
  \item You know the \emph{Nine Measures} (light, draft, echo, dust, taste of iron, sweat-chill, lamp-shadow, bell-lag, head-ache). With a \textbf{Wits + Notice} test (DV 3), detect hidden passages or environmental hazards.
  \item Once/scene, \emph{read a structure like a ledger}: +1 die to understand its construction, weaknesses, or maintenance needs.
\end{itemize}

\subsection*{Cultural Mechanics}

\subsubsection*{Deep Drakes \& Stone-Press (Complication)}
\textbf{Stone-press}—pressure that thinks—warps sums and senses.
\begin{itemize}
  \item \textbf{Fronts:} \emph{Stone-Press [6]}, \emph{Miasma Spread [4]}, \emph{Vent Failure [4]}.
  \item \textbf{In Play:} Failures on perception/planning underground may generate SB that tick these Fronts or impose \emph{Condition} or hallucination tests.
\end{itemize}

\subsubsection*{Tally-Law (Social/Legal)}
\begin{itemize}
  \item \textbf{Oaths:} Breaking a formal oath ticks \emph{Repute −1} and creates \emph{Feud +1}.
  \item \textbf{Boasts \& Sagas:} A public boast creates \emph{Audience: Expectant}; fulfill it for \emph{Repute +1}; fail and mark \emph{Exposure} or \emph{Feud}.
  \item \textbf{Etiquette Hooks (once/scene, Aeler venues):} Present guest-loaf and a lit lantern at a threshold to shift your next social action \emph{one step safer} while you remain a guest. Covering iron or stepping on stone cancels the first SB from rites/negotiations this scene.
\end{itemize}

\subsubsection*{Strings \& Tools}
\begin{itemize}
  \item \textbf{Keystone Tablet}: Establish or pause a route; once/scene, convert a chase into a prepared stand: defenders gain \textbf{Position +1}.
  \item \textbf{Null-Bell}: On ring, cancel one psychic push/compel; costs lamp-time (mark gear wear or Fatigue).
  \item \textbf{Oven Charter Seal}: Force a public bowl in markets—one round of fair-price negotiation before violence may escalate.
  \item \textbf{Sluice Key}: \textbf{DV −1} on operations hinging on water/flood/sanitation; abuse creates \emph{Public Outrage [4]}.
  \item \textbf{Air Scrip}: Negate the first suffocation/miasma consequence underground in a leg; on use, tick \emph{Vent Failure [1]}.
\end{itemize}

\subsubsection*{Display Rights (Status)}
\begin{itemize}
  \item \textbf{Display Charter}: Licensed marks (metal trim, lamp-halos, keystone etching) show public contribution. Fraud draws fines in grain or labor.
  \item \textbf{Rings of Account}: Bands on tools/belts denote kept contracts: water, bread, bridge, mint. Three rings grant first voice in oven courts.
  \item \textbf{Quiet Wealth}: Private hoard without public work is suspect; unworked shine invites auditors.
  \item \textbf{At the table:} Present a valid Display Writ to gain \textbf{Position +1} in civic negotiations once/scene; on a miss, start \emph{Audit Clock [4]}.
\end{itemize}

\subsubsection*{Orders \& Companies (Factions)}
\begin{itemize}
  \item \textbf{Iron Avengers}: Traditionalists who enforce blood-feuds inscribed on keystones.
  \item \textbf{Spirit Shield Warriors}: Ancestor-venerating guards with mask-helms etched in lineage prayers.
  \item \textbf{True Masons}: Wanderers who repair ancient Aeler work.
  \item \textbf{Edgewalkers}: Border scanners who hunt for the profitable gap.
  \item \textbf{Reform Lodges}: Foothill freeholds and city cells arguing for gentler contracts with neighbors.
\end{itemize}
\emph{Example boons/risks:} Mason's Oath, Edgewalker Marker, Gray List Token, Mask Rights Forfeit.

\subsection*{Soft Power: Keystone Diplomacy \& Infrastructure Sovereignty}
Aeler influence travels by \emph{charter, standard, and switch}, not by spear.
\begin{itemize}
  \item \textbf{Standards Bind:} \emph{Mint \& Measure} make markets legible. Cities that adopt Aeler weights gain stable prices—and accept Aeler audit halls. \emph{In play:} presenting stamped measures grants \textbf{DV −1} to enforce contracts or expose fraud.
  \item \textbf{Oven Charters:} \emph{Grain Surety} keeps winters calm. Charter Days obligate fair-pricing courts before force. \emph{In play:} invoke an Oven Charter Seal to require one round of negotiation; cancel the first riot SB this scene.
  \item \textbf{Water Keys:} \emph{Sluice guilds} lease flow, not walls. Trade booms when gates open; siege starves when they close. \emph{In play:} exchanging a Sluice Key with a civic body banks \emph{Public Gratitude [2]}; abuse flips it to \emph{Outrage [4]}.
  \item \textbf{Keystone Clauses:} Bridges and gates include \emph{rate-for-repair}. Default pauses service “politely” until arrears are paid. \emph{In play:} once/scene, declare a keystone pause to impose \textbf{Effect −1} on enemy movement or logistics.
  \item \textbf{Air \& Sanitation:} \emph{Vent courts} and wasteflows curb plague. The city that breathes owes the hand that tuned the vents. \emph{In play:} with plans and access, Aeler gain \textbf{Position +1} to quell disease/panic clocks.
  \item \textbf{Apprentice Exchanges:} Sending masters abroad seeds techniques and loyalties. \emph{In play:} spend a season contact to treat a foreign workshop as Friendly for one operation.
  \item \textbf{Null-Bells \& Audit Halls:} Disciplined speech zones deter panic and glamours. \emph{In play:} ringing a null-bell suppresses one social \emph{fear/panic} tag for an exchange.
\end{itemize}
\paragraph{Soft-Power Clocks}
\begin{itemize}
  \item \emph{Public Gratitude [4]} \(\rightarrow\) discounts, calm crowds, smoother permits.
  \item \emph{Audit Clock [4]} \(\rightarrow\) fines, seizures, reputational scars if you flaunt standards.
  \item \emph{Co-Prosperity [6]} \(\rightarrow\) shared surpluses, joint projects, mutual defense clauses.
\end{itemize}

\subsection*{Suggested Bonds \& Complications}
\paragraph{Bonds}
\begin{itemize}
  \item \textbf{Mason-Brother}: Minor edge on structural assessments/repairs and a contact in a lodge.
  \item \textbf{Vent-Prior's Apprentice}: Guidance in underground survival; access to holds.
  \item \textbf{Oven-Warden's Acquaintance}: Insight into grain courts; leverage in hungry quarters.
  \item \textbf{Sluicewarden's Debt}: A water-math favor owed—potential flood control or redirected trade.
\end{itemize}
\paragraph{Complications}
\begin{itemize}
  \item \textbf{Stone-Press Susceptibility}: In deep stress underground, test \textbf{Spirit} (DV 3) or suffer \(-1\) die to tasks from pressure-sickness.
  \item \textbf{Ledger Dependency}: Separated from records \(>\) 1 day: \textbf{Position −1} on planning/resource management until re-synced.
  \item \textbf{Surety Obligation}: You guaranteed resources to a community. Fail to deliver: mark 2 segments on \emph{Obligation} (or take a lasting Complication).
\end{itemize}

\clearpage
\section{Aelinnel — People of Sums, Bough, and Bright Things}

\subsection*{Background: Gnomes of Stone, Bough, and Bright Things}
The Aelinnel dwell in the hawthorn hills south of the Valewood, their lives intertwined with living wood, worked stone, and precise mathematics. Halls run like veins through granite and thorn; bridges hum when tuned; bargains arrive on two ledgers—what was \emph{said}, and what was \emph{meant}. To walk their country is to feel math underfoot: steps safer when even, doors opening to right sequence, moonlight that prefers tidy logic.

They are fae-kin cousins to the Lethai, smaller in stature, bright-eyed and quick. Society rests on \textbf{Courtesy, Copper, and Count}. Count or be counted. Speak your steps, breaths, or stitches and the land steadies. Favor copper before the courts; copper is polite, naked iron is an insult unless named or gilded. Recite a simple sequence when tension frays to cancel the first misstep in navigation or negotiation.

\paragraph{Law of Sums}
Proceedings in hawthorn courts require three clean courtesies: \emph{no naked iron}; \emph{two-ledger speech} (said \& meant); and \emph{return what points the way} (cords, marks, antlers). Keep these and even thorns hold back; slight them and arches close, time drifts, and messages arrive folded and misaddressed.

\subsection*{Racial Skill Increase}
Choose one:
\begin{enumerate}
  \item \textbf{Copper Courtesy \& Counting Etiquette}: Once/scene, careful counting shifts \textbf{Position +1} for a patterned action (locks, paths, ritual speech). In fae-facing scenes, presenting copper or brass tools negates the first offense penalty.
  \item \textbf{Two-Ledger Talk}: When you clearly state both \emph{said} and \emph{meant}, you may cancel the first social \textbf{SB} against you this scene. If you refuse, the next bargain seeks collateral (memory or name, at the GM's discretion).
\end{enumerate}

\subsection*{Thematic Attribute}
Increase either \textbf{Wits} or \textbf{Spirit} by 1 (to a max of 5). \textbf{Body} and \textbf{Presence} unchanged.

\subsection*{Size \& Equipment Limits}
Aelinnel stature and leverage limit heavy kit.
\begin{itemize}
  \item \textbf{Restriction:} Aelinnel \emph{cannot} wield \textbf{Heavy} weapons or wear \textbf{Heavy} armor. They may freely use \textbf{Light} and \textbf{Medium} weapons/armor suited to their frame.
  \item \textbf{Design Note (at the table):} Attempts to circumvent with makeshift rigs are treated as fiction-only stunts; if allowed, apply \textbf{Position --1} and \textbf{DV +2} and remove any benefits from the \textit{Finesse} tag.
\end{itemize}

\subsection*{Talents}

\subsubsection*{Long Memory (3 XP — Minor)}
\textit{Req: \textbf{Spirit} 2+}
\begin{itemize}
  \item Perfect recall of the last week’s events.
  \item +1 die to \textbf{Lore} and \textbf{Insight} when drawing on historical/cultural detail.
  \item \textbf{Once/session}, surface a crucial long-term detail the GM must render truthfully (scope: a person, place, or clause you directly encountered).
\end{itemize}

\subsubsection*{Cold Reading (3 XP — Minor)}
\textit{Req: \textbf{Wits} 2+, \textit{Long Memory}}
\begin{itemize}
  \item +1 die to \textbf{Sway} and \textbf{Insight} from close observation.
  \item With \textbf{Wits + Insight} (DV 3), infer a motive or pressure the target is concealing.
  \item \textbf{Once/scene}, if you observed quietly for a beat, gain \textbf{Position +1} on your first parley in that scene.
\end{itemize}

\subsection*{Cultural Mechanics}

\subsubsection*{Hawthorn Halls \& Law of Sums}
Paths are counted by antler-posts; tide-cut stairs descend to black sea-rifts; causeys of pale flags show at dawn, at dusk—and whenever someone is counting aloud.
\begin{itemize}
  \item \textbf{Counting Etiquette}: Once/scene, careful count grants \textbf{Position +1} for patterned action.
  \item \textbf{Copper Over Iron}: In fae-facing scenes, copper/brass tools negate the first offense penalty.
  \item \textbf{Two-Ledger Talk}: Stating both \emph{said} and \emph{meant} cancels the first social \textbf{SB} this scene.
  \item \textbf{Hazel Favors}: Earned by restoring way-things (antlers, cords, ferry rights). \textbf{Once/leg}, spend to downgrade a glamour/geas.
\end{itemize}

\subsubsection*{People of Stone \& Bough}
Charcoal-burners read omen by smoke hums; stone-singers soothe walls with low chords; wardens hammer copper nails where iron offends; goat-herds measure danger in hoof-widths.
\begin{itemize}
  \item \textbf{Markets Under Living Roofs}: Weights and measures matter—producing a certified rod cancels the first jurisdiction/commerce snag in that market.
  \item \textbf{Reputation Echoes}: Return way-cords, restore antlers, pay tide-dues to bank \emph{Hazel Favors}.
\end{itemize}

\subsubsection*{Courts, Hunts, and Gates}
Aelinnel powers are etiquette engines and logic traps more than tyrants.
\begin{itemize}
  \item \textbf{Lady of Thorns}: Punishes breaches precisely; rewards perfect sequences.
  \item \textbf{Green Knight}: Duels by paths and proofs, not boasts.
  \item \textbf{Moonlit Ride}: Grants one night’s clemency if you name the right horn-count.
  \item \textbf{Green Gate}: Demands exact change in truths before it opens.
\end{itemize}

\subsubsection*{Tides, Ledgers, and Names}
Rivers and sea-caves carry their own arithmetic.
\begin{itemize}
  \item \textbf{Tide-Reeves}: Filing plans before neap earns a Tide Window; skip the ledger and your next crossing suffers \emph{Wrong Tide}.
  \item \textbf{Seals \& Hours}: Neglected barge seals invite \emph{Wrong Hour}.
  \item \textbf{Said/Meant Receipts}: Some stalls require dual receipts; single-ledger haggling risks a surcharge payable in memory or name.
\end{itemize}

\subsubsection*{Aelinnel Mood: Dark-Wonder}
Paths shorten for those who keep count and lengthen for the proud. Petals fall like knives and settle into proofs; antler-posts rearrange themselves when the land takes offense. Time miscounts, and the sun arrives at the wrong hour with excellent logic.

\subsection*{Strings \& Tools}
\begin{itemize}
  \item \textbf{Antler-Post Bead}: \textbf{Once/scene}, treat a wild path as \emph{Signed}: \textbf{DV −1} to Traverse.
  \item \textbf{Counting Cord}: When stretched and tapped in sequence, grants \textbf{Position +1} on one trap/lock sequence.
  \item \textbf{Copper Nail Kit}: Negate the first \emph{iron-offense} penalty in a scene; if misused, start \emph{Thorn Displeasure [2]}.
  \item \textbf{Two-Ledger Rod}: A certified measure; \textbf{DV −1} to enforce a stall’s weight/price clause.
  \item \textbf{Tide Window Seal}: Mark a safe hour for crossing; \textbf{once/leg}, cancel \emph{Wrong Tide}.
\end{itemize}

\subsection*{Soft Power: Courtesy Clauses \& Gate-Math}
Aelinnel influence travels by etiquette, sequence, and small infrastructure that rewrites how strangers meet.
\begin{itemize}
  \item \textbf{Sequence Rights}: The party that maintains way-things claims first say; \textbf{once/scene}, invoke to require a parley beat before force.
  \item \textbf{Counting Thresholds}: Marked thresholds heed those who count; allies entering on your count gain \textbf{Position +1} on their opener.
  \item \textbf{Receipt Culture}: Dual ledgers make fraud expensive; presenting a said/meant receipt grants \textbf{DV −1} to unwind a trick clause.
\end{itemize}
\emph{Soft-Power Clocks:} \emph{Hazel Favor [4]}, \emph{Thorn Displeasure [4]}, \emph{Right Hour [3]}.

\subsection*{Suggested Bonds \& Complications}
\paragraph{Bonds}
\begin{itemize}
  \item \textbf{Hawthorn Courtier}: Access to fae-kin etiquette and minor favors.
  \item \textbf{Tide-Reeve's Acquaintance}: Tide windows, river law, and ferries.
  \item \textbf{Green Market Broker}: Wrapped truths and two-ledger bargaining.
  \item \textbf{Stone-Singer's Apprentice}: Stone songs and structural lore.
\end{itemize}
\paragraph{Complications}
\begin{itemize}
  \item \textbf{Context-Sensitive Speech}: Reading texts older than two generations requires \textbf{Lore + Notice} (DV 4–5) or a \emph{Context String}. Using archaic registers without keys adds \textbf{DV +1}.
  \item \textbf{Overload Sensitivity}: Too many inputs at once mark \textbf{1 Fatigue} and impose \(-1\) die on next \textbf{Insight/Notice}.
  \item \textbf{Brittle Focus}: The first \textbf{Harm 1 (blunt)} each scene converts to \textbf{1 Fatigue}; resolve further Harm normally.
\end{itemize}

\clearpage
\section{Player's Guide: Aelaerem — People of Hearth \& Hollow}

\subsection*{Background: Halflings of Hearth \& Hollow}
The Aelaerem are a people of movement and assembly, living among gentle slopes and hedged lanes. Small in stature and large in memory, they bind promises with bread and lantern-light, and measure seasons by harvest masks and market bells. Hospitality is their public law; beneath it runs an older hedge-law of cup-marks, red thread, and the quiet attention of the Neighbors. Their ``hearth magic'' is housekeeping writ large: doors set true, lamps trimmed, courtesies kept—less spell than system, a precise regard for seasons, thresholds, and debts.

\paragraph{Hearth-Law \& Guest-Right}
A red door promises bread, salt, and one safe night if you come honest. Entering or hosting ``under bread and lantern'' gentles the next parley; hospitality is both shield and clause.

\subsection*{Racial Skill Increase}
Choose one:
\begin{enumerate}
  \item \textbf{Hearth-Law \& Guest-Right}: +1 die to \textbf{Sway}, \textbf{Lore} (local custom), and \textbf{Survival} (shelter signs). Gain \textbf{Position +1} when properly offering or receiving hospitality.
  \item \textbf{Lantern-Law \& the Wardens}: +1 die to \textbf{Notice}, \textbf{Survival} (paths, omens), and \textbf{Tinker} (lamps, small tools). Gain \textbf{Position +1} when you correctly observe small courtesies or maintenance rites.
\end{enumerate}

\subsection*{Thematic Attribute}
Increase either \textbf{Wits} or \textbf{Presence} by 1 (to a max of 5). \textbf{Body} and \textbf{Spirit} unchanged.

\subsection*{Size \& Equipment Limits}
\begin{itemize}
  \item \textbf{Small-Statured}: Agility and subtlety by build and habit.
  \item \textbf{Restriction}: Aelaerem cannot use \textbf{Heavy Weapons} or \textbf{Heavy Armor}. They excel with \textbf{Light} and \textbf{Medium} kits favoring finesse, mobility, and ward-rituals.
\end{itemize}

\subsection*{Talents}

\subsubsection*{Heightened Senses (3 XP — Minor)}
\textit{Req: \textbf{Wits} 2+}
\begin{itemize}
  \item +1 die to \textbf{Notice} and \textbf{Survival}.
  \item With \textbf{Wits + Notice} (DV 3), detect hidden creatures/objects.
  \item In natural environments, gain \textbf{Position +1} on stealth and tracking.
\end{itemize}

\subsubsection*{Root-Balance (3 XP — Minor)}
\textit{Req: \textbf{Body} 2+, \textit{Heightened Senses}}
\begin{itemize}
  \item +1 die to \textbf{Athletics}; resist shove/knockdown more easily.
  \item Move through natural terrain without penalty.
  \item \textbf{Once/scene}, stabilize on precarious footing with \textbf{Body + Athletics} (DV 3).
\end{itemize}

\subsection*{Cultural Mechanics}

\subsubsection*{Hearth-Law \& Guest-Right}
\begin{itemize}
  \item \textbf{Red Door Hospitality}: Present a guest-loaf token beneath a lit lantern to soften a risky social exchange (\textbf{Position +1}) or cancel the first \emph{strange} complication in a scene.
  \item \textbf{Lantern-Writ}: Simple rites keep bounds sweet—``Bread \& Salt'' (\textbf{Position +1} once/scene), ``Broom Witness'' (establish \emph{Oath [4–6]}), and \emph{Iron-Lace \& Red Thread} (\textbf{Effect +1} vs. compulsion) while properly maintained.
  \item \textbf{The Neighbors}: Leave butter at cup-marks, keep the festival calendar, and count the stiles aloud. Observance smooths the night; neglect invites \emph{Hollow Attention}.
\end{itemize}

\subsubsection*{Lantern-Law \& the Wardens}
\begin{itemize}
  \item \textbf{Count the Load}: Tap a beam three times and listen; \textbf{once/scene}, a measured tap grants \textbf{Position +1} to \textbf{Traverse/Endure} in caves, bridges, or crowded structures.
  \item \textbf{Copper Courtesy}: Copper is polite to stone and honest to labor. Presenting copper tools or a mason's tally negates the first structural environment \textbf{SB}, or grants \textbf{DV −1} to parley with miners/masons/porters.
  \item \textbf{Return the Chalk}: Anything that points the way (chalk, cord, placard) must be restored; doing so cancels the first environment \textbf{SB} this scene.
\end{itemize}

\subsubsection*{The Quiet Powers (Neighbors)}
\begin{itemize}
  \item \textbf{The Pale Shepherd}: \textbf{Once}, by clause and courtesy, a traveler may pass ``uncounted''—unseen by what tallies footfalls under the soil.
  \item \textbf{Hollow Attention}: Breaking hedge-law draws subtle reprisals: bells toll soft, red thread appears where you did not tie it, a door leads briefly elsewhere.
\end{itemize}

\subsubsection*{Seasons of Mask \& Harvest}
\begin{itemize}
  \item \textbf{Mummers}: Keep stricter rules than any priest; the Thresher-King's guard walks in red hoods when fields demand order.
  \item \textbf{Festivals}: For a night, masks legitimate certain crossings; private moots under the Oak settle quarrels; an elder's blessing opens doors that ignore coin.
  \item \textbf{Omens}: Scarecrows watch the lane; lanterns burn blue at the ford; chalk mazes fill with mist; sometimes the Moot Oak bleeds sap the color of wine.
\end{itemize}

\subsubsection*{Trade, Craft, \& Tokens}
\begin{itemize}
  \item \textbf{Keeps}: Cider, perry, beeswax, and wool spin the lane’s economy.
  \item \textbf{Tokens}: Orchard grafts, mill-tokens, shepherd whistles, mover pressings. (Example: a shepherd’s whistle makes dogs and door-bolts heed for one scene.)
  \item \textbf{Hearth Magic}: Red thread binds promises, lantern-writ holds the dark at bay, careful count and courtesy keep thresholds sweet.
\end{itemize}

\subsection*{Strings \& Tools}
\begin{itemize}
  \item \textbf{Guest-Loaf Token}: \textbf{Once/scene}, treat a tense arrival as \emph{Hospitable} (\textbf{Position +1} opener).
  \item \textbf{Lantern Hood}: Convert a bright scene to \emph{Shaded} locally; cancel one glare-based penalty.
  \item \textbf{Red Thread Kit}: \textbf{Once/scene}, bind a simple promise as an \emph{Oath [4]}.
  \item \textbf{Broom Witness}: Establish a household oath; while it ticks, outsiders face \textbf{DV +1} to trespass or deceive under that roof.
\end{itemize}

\subsection*{Suggested Bonds \& Complications}
\paragraph{Bonds}
\begin{itemize}
  \item \textbf{Apple-Matron's Favor}: Hospitality leverage; influence in local markets; feast-clause invocations.
  \item \textbf{Lantern-Warden's Knowledge}: Path marks, safe-passage signs, omen-reading.
  \item \textbf{Mummers' Captain's Acquaintance}: Festival law, mask permissions, night-crossing exceptions.
  \item \textbf{Hedge-Witch's Debt}: Small potent favors—at a cost.
\end{itemize}

\paragraph{Complications}
\begin{itemize}
  \item \textbf{Hollow Stirring [6]}: Courtesies neglected; omens grow frequent and bite.
  \item \textbf{Gloam Choir [6]}: Deeper threat from boundary failures; spiritual or fae pressure escalates.
  \item \textbf{Scale Shock}: Smallness is overlooked in mass conflicts or intimidation; social \textbf{DV +1} to \emph{impose} on much larger foes unless you stand under hospitality or law.
\end{itemize}

\clearpage
\section{Player's Guide: Lethai — Root, River, \& Roof-Tree / Mind’s Eye \& Civic Measure}

\subsection*{Background: Woodwise Lawkeepers \& Civic Engineers}
The Lethai are sundered by an old constraint: no one may bear both the Gift of the Body and the Gift of the Mind. From this division grew two sister cultures:

\paragraph{Lethai-al (Wood-Elves) — Root, River, Roof-Tree}
They dwell where roof-trees braid the sky and rivers think aloud. Their memory is arboreal—rings, seasons, coppice ledgers—and their oaths are \emph{root-law}: debts in years, paid in living work. To outsiders they seem quiet; to neighbors they are relentless auditors of footprint and flow. Strengths: living law, environmental stewardship, contextual craft.

\paragraph{Lethai-thora (City-Elves) — Mind’s Eye, Civic Measure}
They make circles in cities (chiefly Thepyrgos), weighing arguments like bridges and translating other peoples’ law into forms that carry. Their courts count consequences; their speech is context-saturated and exact. Strengths: memory, jurisprudence, logistics, civil design.

\paragraph{First Courtesy}
“Name yourself once; name the river twice; never name the forest as if it were yours.” Lamps are witnesses, paths are clauses, and small laws keep the world sweet.

\subsection*{Racial Skill Increase}
Choose one, keyed to your branch:
\begin{enumerate}
  \item \textbf{Gift of the Body (Lethai-al)}: +1 die to \textbf{Body}-based actions; \textbf{once/scene} you may spend \textbf{1 Boon} to exceed normal limits for one physical action (leap, balance, sprint, scent-track).
  \item \textbf{Gift of the Mind (Lethai-thora)}: +1 die to \textbf{Wits/Spirit}-based actions; \textbf{once/scene} you may spend \textbf{1 Boon} to recall, deduce, or frame context that was beyond immediate knowledge (gloss, precedent, supply path).
\end{enumerate}

\subsection*{Thematic Attribute}
\begin{itemize}
  \item \textbf{Lethai-al}: Increase \textbf{Body} \emph{or} \textbf{Wits} by 1 (max 5).
  \item \textbf{Lethai-thora}: Increase \textbf{Wits} \emph{or} \textbf{Spirit} by 1 (max 5).
\end{itemize}

\subsection*{Talents}

\subsubsection*{Lethai-al Path — Embodied Presence}
\paragraph{Canopy Spring (4 XP — Minor)}
\textit{Req: Heightened Senses; Body 2+}\\
+1 die to climbing, vaulting, branch-run. \textbf{Once/scene} in forest, gain \textbf{Position +1} on a movement action. With \textbf{Body + Athletics} (DV 4), clear a gap others treat as impassable.

\paragraph{Scent of Rain (4 XP — Minor)}
\textit{Req: Heightened Senses; Survival 1+}\\
+1 die to predict weather, smoke, blight. With \textbf{Wits + Survival} (DV 3), sense approaching storm/fire/plague front. Track by scent with \textbf{Effect +1} in natural terrain.

\subsubsection*{Lethai-thora Path — Mental Acuity}
\paragraph{Memory Canticle (4 XP — Minor)}
\textit{Req: Long Memory; Lore 2+}\\
Line-true recall of texts and testimony. +1 die to research/translation. \textbf{Once/scene}, provide crucial context that advances an investigation or reduces a legal/social \textbf{DV} by \textbf{1} if you can cite source and frame.

\paragraph{Number Music (4 XP — Minor)}
\textit{Req: Long Memory; Wits 2+}\\
+1 die to design/repair/logistics. With \textbf{Wits + Craft} (DV 3), solve an engineering/flow proof on the fly. When you \emph{speak the scaffold} (frame steps aloud), related rolls are \textbf{DV −1} this scene.

\subsection*{Cultural Mechanics}

\subsubsection*{The Curse of Division}
\begin{itemize}
  \item \textbf{Rule of Two Gifts}: Choose \emph{Body} (Lethai-al) or \emph{Mind} (Lethai-thora) at creation; take Talents only from that path. Attempting to straddle both imposes \textbf{−1 die} to all rolls until scene end.
  \item \textbf{Season Switch}: Changing paths is a season project with social cost (forfeit one String tied to former gift).
  \item \textbf{Bridge-Born Clause (Rare Half-Elf)}: A quarter-lineage may take \emph{one} Body Gift \emph{and} \emph{one} Mind Gift (see GM for context obligations).
\end{itemize}

\subsubsection*{Shade Etiquette (Lethai-al)}
\begin{itemize}
  \item \textbf{Iron Covered}: Bare iron offends ward-lines; wrap it. First entry with covered iron grants \textbf{Position +1} in parley to pass.
  \item \textbf{Name Once}: Speak name and intent at the edge; do not claim the forest.
  \item \textbf{Step on Stone}: Use laid stones; crushed shoots are debts.
  \item \textbf{Water First}: Pour a first cup for river or cistern before you drink.
  \item \textbf{Leave the Light}: Replant, mend, or pay for shade you take. (See \emph{Light-Dues}.)
\end{itemize}

\subsubsection*{Context-Saturated Speech (Lethai-thora)}
\begin{itemize}
  \item \textbf{Context Keys}: place-name, season-mark, kinship-hand, roof-tree sign. Missing any two invites misreadings.
  \item \textbf{Old Texts}: Manuscripts older than two generations need gloss-trees (marginal twig glyphs) or a songkeeper.
  \item \textbf{In Play}: Reading/pleading in older registers is \textbf{DV +1} unless you hold a \emph{Context String} (gloss-tree, witness, song). Cashing it reduces \textbf{DV −1} and establishes a \emph{Shared Frame} (\textbf{Position +1} opener).
\end{itemize}

\subsubsection*{Strings \& Ledger (Shared)}
\begin{itemize}
  \item \textbf{Light-Dues}: Every fell, ferry, and fire owes a balance in replanting, canal-clearing, or kinder rates.
  \item \textbf{Strings (examples)}: light-due receipt; shade-credit; ferry right; resin share; canoe-lane priority; seed tithe.
  \item \textbf{Use}: Cash a light-due for \textbf{DV −1} on operations framed as repair/replanting/flood-work. Abuse starts \emph{Under-Root Grudge [1]}.
  \item \textbf{Quotas as Clocks}: Boat-Timber Quota [6], Lanternwood Allotment [4]. Fill to unlock export; overfill triggers \emph{Canopy Censure}.
\end{itemize}

\subsubsection*{Patron Ties (Common Lethai Bonds)}
\begin{itemize}
  \item \textbf{Lethai-al}: Often entreat \emph{Inaea} (web, guest-right, line-sanctuary) and \emph{Isoka} (shedding, decisive strike) for rites on path, bridge, and hunt.
  \item \textbf{Lethai-thora}: Favor \emph{The Witness} (truth, record) and \emph{Sacred Geometry} (form, proportion); radicals court the \emph{Clockwork Monad} for “managed process” at moral risk.
\end{itemize}

\subsection*{Strings \& Tools (Table Use)}
\begin{itemize}
  \item \textbf{Shade-Credit}: Treat a contested crossing as \emph{Hospitable}: \textbf{Position +1} opener; on abuse, tick \emph{Under-Root Grudge [1]}.
  \item \textbf{Ferry Right}: \textbf{DV −1} to move crews/loads by water where named.
  \item \textbf{Resin Share}: \textbf{Effect +1} on mend/seal actions; marks \emph{Lanternwood Allotment}.
  \item \textbf{Context String}: Redeem to remove \textbf{DV +1} from archaic law/speech this scene and bank \textbf{+1 Boon} on a clean success.
\end{itemize}

\subsection*{Suggested Bonds \& Complications}
\paragraph{Bonds}
\begin{itemize}
  \item \textbf{Songkeeper’s Trust (al)}: Vouches adherence to shade etiquette; interprets ancient paths and oaths.
  \item \textbf{Warden’s Oath (al)}: Passage rights and woodwise summons when the forest is wronged.
  \item \textbf{Sumwright’s Compact (thora)}: Two-ledger arbitration in mixed courts; contract leverage.
  \item \textbf{Archive-Keeper’s Debt (thora)}: Access to restricted stacks—at a later price in service.
  \item \textbf{Bridge-Born Kinship}: Mutual aid among rare dual-gift lineages navigating both spheres.
\end{itemize}

\paragraph{Complications}
\begin{itemize}
  \item \textbf{Division’s Bite}: Attempting cross-path use (Body \emph{and} Mind gifts in one scene without clause) imposes \textbf{−1 die} to all rolls until scene ends.
  \item \textbf{Context Fragility}: Social/legal \textbf{DV +1} when stripped of context keys; hostile courts may weaponize misframing.
  \item \textbf{Canopy Censure}: Exceeding quotas or slighting root-law triggers censure—\emph{Light-Dues [4]} starts and local passage turns \emph{Controlled}.
\end{itemize}

\clearpage
\section{Player's Guide: Lethai-ar — The Oathbound (Dark Elves)}

\subsection*{Background: The Vowed in Silk \& Scale}
The Lethai-ar are not a separate bloodline so much as a vow-bound cadre. They are Lethai-al (Body-gift) or Lethai-thora (Mind-gift) who step off those paths to live under threshold patrons—\textbf{Inae} the Weaver and \textbf{Isoka} the Serpent. In some ages they are scarce; in others—like the present—they gather wherever borders fray and oaths need teeth. They are known for \emph{mask-right}, precise courtesies, and the unsettling efficiency of vows that bind places, routes, and roles.

\paragraph{Marks Are Common, Discipline Is Distinct.}
Ink, scar, resin-inlay, and rite-born \emph{Marks} exist across many peoples. The Lethai-ar did not invent them and do not own them; they \emph{institutionalize} them—pairing each Mark with context keys (place, season, witness) and ledgered prices, so boons do not drift into curses.

\subsection*{Racial Skill Increase}
Choose one according to patronal leaning:
\begin{enumerate}
  \item \textbf{Weaver’s Reading (Inae)}: +1 die to uncover \emph{connections} (routes, plots, safe-conducts). \textbf{Once/scene}, with \textbf{Wits + Notice} (DV 3), declare one hidden tie (who vouches, what clause binds, which alley joins).
  \item \textbf{Venin Lore (Isoka)}: +1 die to \emph{identify, dose, and remedy} toxins and social “poisons.” \textbf{Once/scene}, convert \textbf{Harm 1 (toxin)} to \textbf{1 Fatigue}.
\end{enumerate}

\subsection*{Thematic Attribute}
Increase \textbf{Wits} \emph{or} \textbf{Spirit} by 1 (max 5). Body and Presence unchanged.

\subsection*{Talents (Entry Paths)}
\paragraph{Needle-Quiet (4 XP — Minor)}
\textit{Req: Stealth 2+}\\
In dim, patterned cover (silk, lattice, shutters), gain \textbf{+1 die} to \textbf{Stealth}. \textbf{Once/scene}, hold still through an exchange as if unseen (you remain targetable only by area or guessed fire).

\paragraph{Fang of Timing (4 XP — Minor)}
\textit{Req: Survival 1+ or Subterfuge 1+}\\
Frame a precise moment (counted breath, bell-beat). Your next action that keys to that beat gains \textbf{Position +1}. \textbf{On a Miss}, mark \textbf{1 Fatigue} (the moment slips).

\subsection*{Cultural Mechanics}

\subsubsection*{Two Courts, One Edge}
\paragraph{Silk Courts of Inae (Pattern \& Mercy)}
Temper: mercy with memory; knots that mend before they bind.\\
Work: multi-party compacts, reweaving custom, sanctuaries under line.\\
Signs: three-strand cords, ledger-ribbons, masks with tear-slits.\\
Law: “Said \& Meant” tied by a visible clause.\\
Sin: binding without consent; repair that erases the harmed.

\paragraph{Coil Courts of Isoka (Change \& Cure)}
Temper: cunning without needless cruelty; a sharp cure offered with the cut.\\
Work: expose weak seams, stage molts (identity exits), pair poison to remedy.\\
Signs: shed-skin sashes, cup-and-vial pairs, scalpels in green thread.\\
Law: every wound must name a reachable remedy.\\
Sin: a wound with no cure; molt forced by shame instead of choice.

\subsubsection*{Oath Etiquette \& Rites}
\begin{itemize}
  \item \textbf{Mask-Right}: Roles declared, then masks donned; speak in role, not over it. Breaking mask-right ends hearing for a season.
  \item \textbf{Speak Twice, Whisper Once}: Say the truth two ways; then whisper the \emph{price} to the witness. If you cannot name the price, you have not made a true ask.
  \item \textbf{Thread Before Blade}: Offer a binding solution first. Steel only where thread was refused.
  \item \textbf{Vial Courtesy}: A dose sits beside its antivenin. Taking one without the other marks bad faith.
\end{itemize}

\subsubsection*{Marks — Shared Craft, Particular Discipline}
Marks are table-facing compacts that grant an edge when \emph{in context} and tilt to curse when unmoored. The Lethai-ar practice is to \emph{key} every Mark to time, witness, and place.

\paragraph{Examples (available to any culture that learns the rite; Lethai-ar formalize keys):}
\begin{itemize}
  \item \textbf{Spider-Bride Mark (Inae)}: \emph{Gift}: \textbf{Position +1} when protecting named guests “under your line”; line-sanctuary can be declared once/scene. \emph{Keys}: guest-roll, lifted lamp, knot-book. \emph{Curse (keyless)}: hospitality fixation—\textbf{DV +1} to withdraw protection even when prudent.
  \item \textbf{Widow’s Spool (Inae)}: \emph{Gift}: lay a hair-fine traversal/trap line once/scene (treat as \emph{Trap [2]} or safe step). \emph{Keys}: knot register, bell-note. \emph{Curse}: path hunger—compelled to “finish the old route” at bad moments.
  \item \textbf{First Shedding Mark (Isoka)}: \emph{Gift}: escape a label/bond once/scene (\emph{Shed-Skin Escape}); clear \textbf{1 Fatigue} on successful exit. \emph{Keys}: warm draught, ash-milk, witness to new name. \emph{Curse}: identity chill—\textbf{Position −1} in cold or when unnamed.
  \item \textbf{Forked Sight (Isoka)}: \emph{Gift}: when a hard truth is spoken in the exchange, your action gains \textbf{Effect +1}. \emph{Keys}: bell-pattern, truth-token. \emph{Curse}: social sting—on Miss, mark \textbf{1 Fatigue} and start \emph{Bruised Pride [2]}.
\end{itemize}

\subsubsection*{War Without Battle (Doctrine)}
Lethai-ar end campaigns by route, ledger, and night.
\begin{itemize}
  \item \textbf{Hedge War}: Re-knit ward-lines to channel intruders into dead ground watched by wardens.
  \item \textbf{River Denial}: Ferries “rest,” weirs open at dusk, mills idle to silt the only footing.
  \item \textbf{Night Lanes}: Silk trip-lines and warning strings; one strike, then silence.
  \item \textbf{Canopy Runners}: Move above sight-lines; arrows fall where footfalls never were.
\end{itemize}

\subsubsection*{Neighbors \& Borders}
\begin{itemize}
  \item \textbf{Ykrul}: Ritual distance, mutual measure. Bowl \& Board seals routes; Ykrul price crossings, Lethai-ar price behavior.
  \item \textbf{Aeler}: Oath-friction underground—lamp-law vs mask-right. Both enforce receipts; argue which witness counts.
  \item \textbf{Lethai-al / Lethai-thora}: The ar recruit from the al; they debate context with the thora. Kinship slows quarrels; thresholds decide them.
\end{itemize}

\subsection*{Strings \& Tools (Table Use)}
\begin{itemize}
  \item \textbf{Bride-Line Writ}: Declare a \emph{protected path} for one leg; trespassers treat the route as \emph{Controlled} and suffer \textbf{DV +1} to press.
  \item \textbf{Vial Pair}: Carry dose \& cure; \textbf{DV −1} to treat poison, and parley \textbf{Position +1} when you show both.
  \item \textbf{Knot-Book}: Spend to assert a clause remembered “in the tying”: \textbf{DV −1} to enforce or unwind an oath once/scene.
  \item \textbf{Mask Ledger}: Track roles \& prices; \textbf{+1 die} to \textbf{Sway} when all parties wear declared masks.
\end{itemize}

\subsection*{Clocks \& Fronts}
\begin{itemize}
  \item \textbf{Mask Integrity [4]} (roles fray); \textbf{Guest-Right Strain [4]} (hospitality abused); \textbf{Coil Paranoia [4]} (cure withheld); \textbf{Web Ossifies [4]} (mercy becomes trap).
  \item \textbf{SB Menu}: \emph{Mask Slips} (role confusion), \emph{Knot Bites} (unpaid clause triggers), \emph{Cold Hour} (Position −1 unless warmed), \emph{Witness Arrives} (price must be named now).
\end{itemize}

\subsection*{Play Hooks}
\begin{enumerate}
  \item \textbf{The Unpriced Mercy}: A sanctuary knot holds a murderer. Name a price that mends without erasing the harmed—under three masks—before dawn.
  \item \textbf{Molt for a City}: A tyrant’s captain begs a molt. The cure exists; its price may break the garrison’s oath-chain.
  \item \textbf{Thread Across Pasture}: A silence-furlong through border pasture is broken; wolves and wardens close. Re-stitch or accept a biting line.
\end{enumerate}

\subsection*{Suggested Bonds \& Complications}
\paragraph{Bonds}
\begin{itemize}
  \item \textbf{Fellow Oath-Bearer (Inae/Isoka)}: Shared vows, shared leverage.
  \item \textbf{Mask-Maker}: Teaches form; vouches mask-right in foreign courts.
  \item \textbf{Patron’s Herald}: Carries omens and small dispensations.
  \item \textbf{Kin in al/thora}: Bridge to broader Lethai custom and correction.
  \item \textbf{Neutral Arbiter (Sumwright/Archive-Keeper)}: Mediates Said/Meant disputes.
\end{itemize}

\paragraph{Complications}
\begin{itemize}
  \item \textbf{Oath-Breaker’s Stigma}: Mask-right denied until penance.
  \item \textbf{Mark Adrift}: A beloved Mark lost its last context—acts as a curse until re-keyed.
  \item \textbf{Mask vs Lamp}: Whose witness rules—silk or single lamp?
  \item \textbf{Patron’s Disfavor}: Omen of chill silk or dry scales; expect a demanded price.
\end{itemize}

\subsection*{Using Lethai-ar at the Table}
Lean on \emph{roles, routes, and remedies}. Put \textbf{price} on every protection, and \textbf{context keys} on every edge—so that when pressure comes, the choice to pay or to cut is clean, witnessed, and costly in a way the table can feel.
\clearpage
\section{Player's Guide: Ykrul — The People of the Violet Steppe}

\subsection*{Background: The Violet Steppe}
The Ykrul are a fierce, pragmatic people of the \textbf{Violet Steppe}. Outsiders often call them \emph{orcs}—a slur that flattens a rich culture into stereotype. Among those who deal fairly with them, they are \textbf{Ykrul}: “the People of the Violet Steppe.”

Their shared memory begins with the \textbf{Great Wake}—a legendary flight from a wrong sky, crossing an ocean by boat and star. From that passage they learned to read \textbf{flow} (what moves), \textbf{weight} (what holds), and \textbf{exit} (what opens). These principles, codified in the sacred geometry of \textbf{Kon'reh}, inform their warfare, logistics, law, diplomacy, and daily life. They bank on a fearsome name on the open plain—and win by \textbf{routes, supply, and parleys} that make blades unnecessary.

Ykrul life is plural: herders, pilots, weavers, judges, scouts, factors, captains. Warfare is \emph{a tool}, not an identity. Their ethics are spatial: to \emph{place well} is to \emph{behave well}—no route you cannot defend, no exit you build that closes behind weaker feet.

Institutions range from \textbf{Meadow Judges} who roam with bowl and board, to \textbf{Kon'reh Masters} who arbitrate by geometry, to \textbf{Stone-Sons \& Rope-Daughters} who prove passages in mountain night, and \textbf{Wake-Wrought} families who list their hulls like saints. Their colors favor violet (grass), gray (stone), blue-white (sky/wake), and motifs of rings, crosses, and stepped lines.

\subsection*{Racial Skill Increase}
Choose one:
\begin{enumerate}
  \item \textbf{Flow/Weight/Exit (Plainscraft)}: +1 die to \textbf{Survival}, \textbf{Tactics}, and \textbf{Lore} about movement, terrain, logistics, and “reading” a situation for viable paths or pressure points.
  \item \textbf{Kon'reh Logic (Sacred Geometry)}: +1 die to \textbf{Insight} and \textbf{Command} when negotiating, planning, or reasoning spatially. \textbf{Once/scene}, you may treat a failed negotiation as a \emph{Partial} if you reframe it with a clean geometric metaphor or map.
\end{enumerate}

\subsection*{Thematic Attribute}
Increase \textbf{Body} \emph{or} \textbf{Wits} by 1 (max 5). Spirit and Presence unchanged.

\subsection*{Talents}
\paragraph{Herd-Mastery (3 XP — Minor)}
\textit{Req: Survival 1+, Presence 2+}\\
+1 die to herding/animal handling. Calm a panicked beast or steer a herd through bad ground with \textbf{Presence + Survival} (DV 3). \textbf{Once/scene}, read herd motion for advantage (\textbf{+1 die} to a linked \textbf{Notice} or \textbf{Tactics} test).

\paragraph{Weather-Reading (4 XP — Minor)}
\textit{Req: Wits 2+, Survival 1+}\\
+1 die to predict weather or navigate it. Sense major shifts hours early with \textbf{Wits + Survival} (DV 3). In open country, acting on a correct forecast grants \textbf{Position +1} to travel/survival actions.

\subsection*{Cultural Mechanics}

\subsubsection*{Wake-Law: The Great Migration}
\begin{itemize}
  \item \textbf{Survival through Adaptation}: Read \emph{flow, weight, exit} before you commit.
  \item \textbf{Kon'reh (Sacred Geometry)}: Meadow (pressure \& path), River (flow \& change), Stone (weight \& witness). Strategy is ethic; a clean placement is a clean duty.
  \item \textbf{Wake Names}: A second name granted by the sea—callable once for an unlooked-for courtesy “as kin on the crossing.”
  \item \textbf{Salt Line}: A coil of salt-stiff rope uncoiled before grave talk—placing all present under Wake-law: plain words, straight debts, no riddles.
\end{itemize}

\subsubsection*{Law, Diplomacy, War—Seen in Lines}
\begin{itemize}
  \item \textbf{Guest Right}: A guest cup at the outer fire. Theft “under smoke” brands your tent-marks gray for a season.
  \item \textbf{Two Ledgers}: Said \emph{and} Meant recorded together; offering both averts face-traps. Lying on \emph{Said} is a grave offense.
  \item \textbf{Blood-Price}: Paid in animals, salt, or length of rope; refusal enters \textbf{Red Weather} (others treat you as a walking storm).
  \item \textbf{Silence Furlong (with Lethai-al)}: A speechless border—no grazing, no felling, no names. Cross in silence, then speak once. Kept, it warms councils; broken, gray-fletched messengers arrive at dusk.
  \item \textbf{Pass \& Harbor Doctrine}: In mountains, hold what stone will bear and promise no more. At sea, get there first \emph{or sing the storm together}; a shared song outranks a sharp keel.
\end{itemize}

\subsubsection*{People \& Institutions}
\begin{itemize}
  \item \textbf{Meadow Judges}: A traveling trio with bowl, board, and braid; their ruling holds one season and one road.
  \item \textbf{Kon'reh Masters}: Geometry arbiters and teachers; respected even by rivals; will play anyone who brings a decent board and reason.
  \item \textbf{Stone-Sons \& Rope-Daughters}: Mountain orders proving night crossings; their braids anchor Aeler engines and make Aeler officers polite.
  \item \textbf{Wake-Wrought}: Sea families who list hulls like saints; name a stolen boat’s rivets and you can claim it under Wake-law.
\end{itemize}

\subsubsection*{Ykrul Ways (Four Grounds)}
\begin{description}
  \item[Meadow Commons] Violet grass, ring-camps, shallow lakes. Fast musters, moving markets, distance diplomacy. Gifts: Herd-Mastery; Weather-Reading; Caravan Craft; Route-Planning (Kon'reh framing).
  \item[Mountain Holds] Knife ridges, pass-stones. Holding lines, winter stores, signal sense. Gifts: Stone-Sense; Counterweight Engineering; Rope-Craft; Avalanche Reading.
  \item[Salt Coasts] Rocky inlets, island runs, river mouths. Pilots, moots, storm windows. Gifts: Storm-Seamanship; Harbor Dues; Shoal Mapping; Blue Moot Etiquette.
  \item[Eastern Steppes] Sky cairns, ward-storms. Exit-finding, omen-reading, cross-cultural guides. Gifts: Cairn-Talking; Ward-Storm Guidance; Long-Leg Logistics; Silence Furlong Etiquette.
\end{description}

\subsubsection*{How Ykrul Win (Beyond the Blade)}
\begin{itemize}
  \item \textbf{Banked Fear (Reputation Economy)}: On entering a venue where your fierce name precedes you, mark \emph{Banked Fear (1)}. \emph{Spend 1}: force \emph{Parley First} (one roll of talk before a fight) \emph{or} shift one enemy action to \emph{Controlled} (they flinch). Bluff and fail to deliver? Erase all Banked Fear until you win publicly.
  \item \textbf{Logistics Edge (Strings \textrightarrow\ DV)}: Convert a cache String (fodder lot, water right, hidden wharf) into \textbf{DV −1} on a linked leg/score. Starve innocents with it and flip the benefit: \textbf{DV −1} becomes \emph{Public Debts +1}.
  \item \textbf{Kon'reh Arbitration (Geometry of Mercy)}: Model roads/exits/lanes in talk. On success, create a \emph{Seasonal Concession} String at a ford/harbor/pass; either side may call it \emph{once/season} without offense.
\end{itemize}

\subsection*{Strings \& Tools (Table Use)}
\begin{itemize}
  \item \textbf{Salt Line Rope}: Uncoil to place a scene under Wake-law; \textbf{Position +1} to resolve disputes cleanly once/scene.
  \item \textbf{Wake-Name Token}: Cash for an unlooked-for courtesy—\textbf{DV −1} to a border, muster, or harbor ask.
  \item \textbf{Meadow Judge’s Braid}: Present to shift a brawl into \emph{Bowl \& Board} arbitration; opens a \emph{Concession [4]} clock both sides can tick.
  \item \textbf{Kon'reh Board \& Stones}: Lay it out to reframe a negotiation as placement; on a clean success, bank \emph{Route Clause} (once/leg: \textbf{DV −1} to movement/supply).
  \item \textbf{Violet Standard}: Raise to claim right-of-parley for your band; first hostile SB becomes \emph{Muttered Threats} (no immediate violence).
\end{itemize}

\subsection*{Clocks \& Fronts}
\begin{itemize}
  \item \textbf{Red Weather [6]} (unpaid blood-price shadows your camp)
  \item \textbf{Feud Ignites [4]} (a slight tends toward blood)
  \item \textbf{Silence Furlong Breach [4]} (border rites strained)
  \item \textbf{Supply Drag [4]} (overextended routes sap will)
  \item \textbf{Gray Marks [4]} (tent-brand shame limits hospitality)
\end{itemize}
\textbf{SB Menu (Steppe):} Dust Line (visibility warps), Dry Kettle (water tighter), False Ford (route misread), Storm Edge (forecast arrives early), Horse Nerves (mounts spook).

\subsection*{Play Hooks}
\begin{enumerate}
  \item \textbf{The Board at Dusk}: Two caravans claim the same ford. Lay the Kon'reh board, win the concession without drawing steel, or face \emph{Red Weather}.
  \item \textbf{Salt on the Wind}: A Wake-Wrought hull is stolen; name its rivets under Wake-law and seize it back mid-moot—without lighting the harbor to war.
  \item \textbf{The Gray Tent}: A guest stole under smoke; your marks run gray. Pay the price in rope, stock, or service before neighboring bands treat you like a storm.
  \item \textbf{The Silent Strip}: A Silence Furlong was trampled during a hunt. Repair rites with the Lethai-al before gray-fletched arrows and cold courtesy freeze trade.
\end{enumerate}

\subsection*{Suggested Bonds \& Complications}
\paragraph{Bonds}
\begin{itemize}
  \item \textbf{Meadow Judge Acquaintance}: A traveling arbiter who can seat disputes and grant rulings that hold for a season.
  \item \textbf{Kon'reh Master’s Respect}: Earned by fair play or clean placement; grants edge when a scene is “set as Board.”
  \item \textbf{Stone-Son/Rope-Daughter Initiate}: Mountain-tested; access to passes, counterweight tricks, and Aeler goodwill.
  \item \textbf{Wake-Wrought Kin}: Sea-law, pilot lore, and a harbor that remembers your name.
  \item \textbf{Foster-Bond (Ykrul \textrightarrow\ Vilikari)}: A formal exchange that opens trade routes and softens borders.
\end{itemize}

\paragraph{Complications}
\begin{itemize}
  \item \textbf{Red Weather}: Your band owes blood-price; hospitality chills until it’s paid.
  \item \textbf{Feud Ignites [4]}: A live quarrel trends toward blood unless cooled by price, proof, or play.
  \item \textbf{Silence Furlong Breach}: Border rites violated; messengers with gray fletchings are on the way.
  \item \textbf{Oath-Breaker’s Shame}: Guest-Right or Said/Meant betrayed; trust collapses across rings.
  \item \textbf{Kon'reh Misstep}: You promised an exit you cannot defend. Publicly re-learn (demonstrate competence) or lose face.
  \item \textbf{Banked Fear Debt}: You spent the name and failed to deliver—erase Banked Fear and invite challengers.
\end{itemize}

\section{Player's Guide: Mixed Heritage — Half-Elves, Half-Ykrul, Half-Others}

\subsection*{Background: Children of Crossings}
The lands of Fate’s Edge are broad and braided. People travel, trade, swear oaths, and fall in love across borders. From these crossings come folk of \textbf{mixed heritage}. You might be the child of a Lethai merchant and a Vilikari factor, a half-Ykrul born in a ford-town to a human parent, or the grandchild of an Aeler vent-prior and a Valewood wanderer. Your identity is not a single stamp but a weave of places, customs, and kin.

Being of mixed heritage is seldom simple. Welcome in one court, weighed in another; fluent in two etiquettes, fully at home in neither. Some will see a bridge, others a trespass. The shape you keep is yours to claim.

This guide favors \emph{reflavoring} the existing race frameworks rather than inventing wholly new rule blocks, keeping focus on narrative flexibility over assumptions about biology.

\subsection*{Creating a Mixed Heritage Character}
\begin{enumerate}
  \item \textbf{Choose One Core Racial Package.} Select one background (Aeler, Aelinnel, Aelaerem, Lethai-al/\,-thora/\,-ar, Ykrul) as your \emph{primary} cultural foundation. This sets your \textbf{Racial Skill Increase}, \textbf{Thematic Attribute}, \textbf{Talents} access, and \textbf{Cultural Mechanics}.
  \item \textbf{Reflavor One Element.} Take one talent, skill bonus, or cultural mechanic from a \emph{second} culture and \textbf{reflavor} it to fit your mixed upbringing. It must make sense in your backstory and present fiction.
\end{enumerate}

\subsection*{Reflavoring Examples (Guidance, Not Limits)}
\paragraph{Half-Ykrul / Half-Lethai-al}
\emph{Herd-Mastery (Ykrul)} reads people as a “crowd-herd” or tracks wildlife lanes through canopy: \textit{“+1 die to understand group motion, crowd dynamics, or animal patterns in your home terrain.”}

\paragraph{Half-Aelinnel / Half-Human (Vilikari)}
\emph{Number Music (Aelinnel)} becomes market math: \textit{“Perform complex trade/logistics calculations with \textbf{Wits+Craft} (or Streetwise) at DV 3; frame deals in clean sums for \textbf{DV −1} once/scene.”}

\paragraph{Half-Aelaerem / Half-Ykrul}
\emph{Kon'reh Logic (Ykrul)} as pantry-sense or hall-placement: \textit{“Once/scene, treat a failed social/negotiation roll as Partial if you reframe with spatial/structural metaphor that fits the venue.”}

\paragraph{Half-Lethai-ar (Inae) / Half-Aelaerem}
\emph{Stillness (Lethai-ar)} becomes host-invisibility: \textit{“+1 die to Stealth in domestic/social bustle; once/scene blend into service and go ‘unnoticed’ for one exchange.”}

\paragraph{Half-Lethai-thora / Half-Vilikari}
\emph{Two-Ledger Talk (Aelinnel)} broadens: \textit{“State both public stance and likely hidden price to cancel the first social SB in a negotiation scene.”}

\paragraph{Half-Ykrul / Half-Linns}
\emph{Storm-Seamanship (Ykrul Coasts)} applies on land: \textit{“+1 die to forecast weather/navigate adverse conditions on steppe or water; on a correct call, gain \textbf{Position +1} for a travel/survival action.”}

\subsection*{Racial Skill Increase (Mixed Heritage)}
Choose one:
\begin{enumerate}
  \item \textbf{Adaptive Skills.} +1 die to \emph{two} different skills that reflect your blended upbringing (e.g., \textbf{Sway} + \textbf{Survival}, or \textbf{Craft} + \textbf{Notice}).
  \item \textbf{Cultural Synthesis.} Take the Racial Skill Increase from your chosen Core Package, and justify how it expresses both sides of your background in play.
\end{enumerate}

\subsection*{Thematic Attribute (Mixed Heritage)}
Choose the Thematic Attribute from your Core Package (e.g., \textbf{Wits} for Lethai-thora, \textbf{Body} for Ykrul) and anchor it in your mixed story (“sharp Wits from city schooling, tempered by border pragmatism”).

\subsection*{Talent (Mixed Heritage)}
Select a Talent from your Core Package \emph{or} one of the flexible options below, and reflavor its fiction to fit your synthesis.

\subsubsection*{Border Walker’s Instinct (4 XP — Minor)}
\textit{Req: Presence 2+, Survival 1+}\\
+1 die to navigate customs, rites, and cross-cultural norms. Sense a border mood with \textbf{Wits+Notice} (DV 3). \textbf{Once/scene}, gain \textbf{Position +1} when mediating between groups or leveraging one culture’s etiquette within another’s court.

\subsubsection*{Tongue of Many Waters (3 XP — Minor)}
\textit{Req: Wits 2+, Sway 1+}\\
+1 die to \textbf{Sway} and \textbf{Insight} with unfamiliar dialects/backgrounds. \textbf{Once/session}, establish a basic pidgin/gesture channel with those who share no tongue (enough for trade or parley). \textbf{Once/scene}, gain +1 on one social roll when you correctly cite a custom from a culture \emph{not} dominant in the interaction.

\subsection*{Cultural Mechanics (Mixed Heritage)}
\begin{itemize}
  \item \textbf{Hybrid Customs.} You may observe simplified forms of two cultures’ mechanics (e.g., Guest-Right tokens \emph{and} Two-Ledger receipts) with narrower scope unless adopted by a host community.
  \item \textbf{Bridging Role.} Treat mixed heritage as a standing fiction tag; some scenes start \textbf{Position +1} for mediation, others \textbf{Position −1} where purity is prized. Let the table lean into both.
  \item \textbf{Reflected Mechanic.} You can mirror the \emph{effect} of a cultural currency (e.g., Reputation/Banked Fear) via deeds across both sides, even if the name differs.
\end{itemize}

\subsection*{Suggested Bonds \& Complications}
\paragraph{Bonds}
\begin{itemize}
  \item \textbf{Family Ties Across Borders.} Kin in two cultures offer haven, rumors, or leverage.
  \item \textbf{Cultural Mentor.} A teacher who drilled you in a rite, register, or craft from one side.
  \item \textbf{Found Family.} A circle of other mixed or liminal folk who vouch when lineage won’t.
  \item \textbf{The Mediator.} Known for seating quarrels between your parent communities.
\end{itemize}

\paragraph{Complications}
\begin{itemize}
  \item \textbf{Identity Fray.} Moments of hesitation or self-editing under scrutiny; first social Miss in a formal venue starts \emph{Doubt [2]}.
  \item \textbf{Dual Expectations.} Two elders pull you in opposite directions; clocks compete (\emph{Obligation A} vs \emph{Obligation B}).
  \item \textbf{Stereotype or Suspicion.} Purists tick \emph{Exposure} on contact; proof is demanded twice.
  \item \textbf{Lost Inheritance.} A rite or ledger you never received; unlock by quest or sponsorship.
  \item \textbf{Translator’s Burden.} Always asked to explain; \emph{once/scene} you may turn that burden into leverage (Position +1) if you take on a new minor obligation.
\end{itemize}
